% !TEX program = xelatex
% =====================================================================
%  BÁO CÁO ĐỒ ÁN TRÍ TUỆ NHÂN TẠO, HỌC PHẦN AIP202
%  Đề tài: FAMS — Hệ thống quản lý sinh viên tích hợp điểm danh bằng
%          nhận diện khuôn mặt
%  Khoa Công nghệ Thông tin, Trường Đại học Kiến trúc Đà Nẵng
%
%  Build:  latexmk -xelatex do-an.tex
%  Overleaf: Menu > Compiler > XeLaTeX, Main document > do-an.tex
%
%  Khung tài liệu dùng lại mẫu đồ án AIP202 của Khoa CNTT (styles/dau.sty,
%  styles/settings.tex giữ nguyên, KHÔNG sửa). Nội dung là của đồ án FAMS.
%
%  SỐ LIỆU TỰ ĐỘNG: file auto/du-lieu-du-an.tex do tools/quet-du-an.py
%  sinh ra từ mã nguồn thật. Không sửa tay file đó, chạy lại script khi
%  dự án thay đổi (xem README.md trong thư mục này).
% =====================================================================
\documentclass[12pt, a4paper, oneside]{book}
% !TEX root = ../do-an.tex
% =====================================================================
%  styles/settings.tex
%  Các thiết lập chung, nạp từ file chính do-an.tex.
%  SINH VIÊN KHÔNG SỬA FILE NÀY.
%
%  Lưu ý: lệnh khai báo lớp tài liệu nằm ở file do-an.tex chứ không nằm
%  ở đây. Lý do: Overleaf dò file chính bằng cách tìm chuỗi khai báo lớp
%  tài liệu trong các file. Nếu chuỗi đó xuất hiện ở file này, Overleaf
%  sẽ biên dịch nhầm file này và báo lỗi "no legal end found", vì file
%  này chỉ chứa thiết lập chứ không chứa phần thân tài liệu.
% =====================================================================

% Toàn bộ nhận diện DAU: phông, màu, bìa, hộp trình bày, header
\usepackage{styles/dau}
\dauLayout

% Giãn dòng và khoảng cách đoạn theo quy cách đồ án
\linespread{1.25}
\usepackage[skip=8pt plus 1pt, indent=15pt]{parskip}

% Ghi chú: mẫu gốc dùng gói atbegshi để vứt trang đầu, vì trang bìa của
% mẫu gốc dựng bằng minipage nên sinh ra một trang trắng. Trang bìa ở
% đây dựng bằng môi trường titlepage nên không có trang trắng, và nếu
% giữ thủ thuật đó thì chính trang bìa sẽ bị vứt mất.

% Bang sinh tu dong co the dai hon mot trang: xltabular = tabularx tu ngat trang.
\usepackage{xltabular}

% Số liệu quét từ kho mã nguồn (số bảng, số endpoint, số test, ngưỡng
% cấu hình, mã commit...). Sinh lại bằng: python tools/quet-du-an.py
% File nay do baocao/tools/quet-du-an.py sinh ra. KHONG SUA TAY.
% Lan quet gan nhat: 12/11/2026 tren commit b099462.
% Sinh lai:  python baocao/tools/quet-du-an.py

\newcommand{\FsoEndpoint}{100}
\newcommand{\FsoModuleApi}{16}
\newcommand{\FsoBang}{21}
\newcommand{\FsoMigration}{15}
\newcommand{\FsoVaiTro}{9}
\newcommand{\FsoTest}{278}
\newcommand{\FsoFileTest}{27}
\newcommand{\FsoManHinh}{34}
\newcommand{\FsoNhomManHinh}{8}
\newcommand{\FsoDongPython}{6.649}
\newcommand{\FsoDongDart}{11.190}
\newcommand{\FsoDongTest}{4.096}
\newcommand{\FcommitCuoi}{b099462}
\newcommand{\FnhanhGit}{feature/flutter-phan-quyen}
\newcommand{\FsoCommit}{31}
\newcommand{\FngayCommit}{05/11/2026}
\newcommand{\FngayQuet}{12/11/2026}
\newcommand{\FchieuVector}{512}
\newcommand{\FnguongNhanDien}{0{,}45}
\newcommand{\FnguongTrung}{0{,}6}
\newcommand{\FnguongThat}{0{,}5}
\newcommand{\FsoKhungXacNhan}{2}
\newcommand{\FphutMuon}{15}
\newcommand{\FphutMoSom}{15}
\newcommand{\FphutAnHan}{30}
\newcommand{\FngayLuuAnh}{90}
\newcommand{\FphutHetHanUrl}{60}
\newcommand{\FphutToken}{480}
\newcommand{\FngayRefresh}{30}
\newcommand{\FsaiSoGps}{100}
\newcommand{\FphutOtp}{10}
\newcommand{\FsoLanSaiToiDa}{5}
\newcommand{\FbanKinhPhong}{50}
\newcommand{\FgiayDoiQr}{30}
\newcommand{\FsoLanGuiEmail}{3}
\newcommand{\FphienBanDongY}{2026-10}
\newcommand{\FtiLeCanhBao}{20}


% Chỗ trống chờ số liệu đo thật hoặc ảnh chụp màn hình. In đỏ đậm để không
% thể nộp nhầm. Tìm chuỗi "CẦN BỔ SUNG" trong PDF trước khi in quyển.
\newcommand{\cando}[1]{{\color{daumaroon}\bfseries[CẦN BỔ SUNG: #1]}}

% ---------------------------------------------------------------------
%  Các bảng sinh tự động: nạp ở đây thành lệnh, thân bảng chỉ gọi tên lệnh.
%  Không đặt \input bên trong tabularx, vì tabularx quét thân bảng nhiều
%  lượt nên \input ở trong sẽ làm vỡ căn lề ("Misplaced \noalign").
% ---------------------------------------------------------------------
% Sinh tu doi tuong APIRouter trong backend/app/api/. KHONG SUA TAY.
\newcommand{\BangApi}{%
\code{/audit-logs} & 1 & Nhật ký hệ thống: ai thao tác gì, lúc nào, thay đổi những gì \\
\code{/auth} & 7 & Đăng nhập, làm mới token, đổi và quên mật khẩu, thông tin tài khoản \\
\code{/classes} & 12 & Lớp học phần, ghi danh, lịch học của lớp \\
\code{/faces} & 6 & Đăng ký khuôn mặt, duyệt hoặc từ chối, dữ liệu khuôn mặt của sinh viên \\
\code{/lecturers} & 6 & Hồ sơ giảng viên, lịch dạy trong ngày \\
\code{/locations} & 4 & Địa điểm điểm danh và bán kính cho phép khi tự điểm danh bằng GPS \\
\code{/evidence} & 1 & Phát ảnh bằng chứng qua đường dẫn được ký điện tử (HMAC), có hạn dùng \\
\code{/notifications} & 4 & Hộp thư thông báo điểm danh của phụ huynh \\
\code{/parents} & 12 & Tài khoản phụ huynh, liên kết với sinh viên, tuỳ chọn nhận thông báo \\
\code{/reports} & 5 & Báo cáo chuyên cần, cảnh báo vắng, xuất Excel, thống kê tổng hợp \\
\code{/rooms} & 5 & Danh mục phòng học: toà nhà, tầng, sức chứa, toạ độ riêng của phòng \\
\code{/schedules} & 2 & Tra cứu lịch học và thời khoá biểu theo tuần \\
\code{/sessions} & 14 & Phiên điểm danh: mở, quét camera, mã QR xoay vòng, tự điểm danh, xác nhận tay, đóng phiên \\
\code{/staff} & 5 & Tài khoản cán bộ phòng ban, gồm Văn phòng Khoa và Lãnh đạo Khoa \\
\code{/students} & 12 & Hồ sơ sinh viên, đồng ý sinh trắc học có phiên bản và rút đồng ý, chuyên cần cá nhân \\
\code{/units} & 4 & Đơn vị trong trường: Khoa, Phòng Đào tạo, Phòng CTSV, Ban ngành \\
}
            % \BangApi
% Sinh tu backend/app/db/models.py. KHONG SUA TAY.
\newcommand{\BangDuLieu}{%
\code{users} & Tài khoản đăng nhập dùng chung cho mọi vai trò \\
\code{units} & Đơn vị: Khoa, Phòng Đào tạo, Phòng CTSV, Ban ngành \\
\code{staff} & Hồ sơ cán bộ, mỗi cán bộ thuộc đúng một đơn vị \\
\code{students} & Hồ sơ sinh viên; đồng ý sinh trắc: phiên bản, cách ghi nhận, người ghi nhận \\
\code{lecturers} & Hồ sơ giảng viên \\
\code{classes} & Lớp học phần trong một học kỳ \\
\code{enrollments} & Ghi danh sinh viên vào lớp học phần \\
\code{locations} & Địa điểm điểm danh kèm bán kính cho phép \\
\code{rooms} & Phòng học thuộc một địa điểm, có thể có toạ độ và bán kính riêng \\
\code{schedules} & Lịch học cố định theo thứ trong tuần \\
\code{attendance\_sessions} & Phiên điểm danh của một lịch học trong một ngày \\
\code{attendance\_records} & Bản ghi điểm danh của từng sinh viên trong phiên \\
\code{attendance\_audits} & Nhật ký mỗi lần sửa trạng thái điểm danh \\
\code{face\_embeddings} & Vector đặc trưng khuôn mặt 512 chiều và trạng thái duyệt \\
\code{revoked\_tokens} & Access token đã bị thu hồi trước hạn \\
\code{refresh\_tokens} & Refresh token của app mobile, lưu bản băm và xoay vòng \\
\code{password\_resets} & Yêu cầu đặt lại mật khẩu bằng OTP \\
\code{parents} & Tài khoản phụ huynh \\
\code{parent\_students} & Liên kết phụ huynh với sinh viên \\
\code{notifications} & Thông báo điểm danh gửi phụ huynh, kèm trạng thái gửi email \\
\code{audit\_logs} & Nhật ký hệ thống: người thao tác, hành động, đối tượng, Khoa liên quan \\
}
   % \BangDuLieu
% Sinh tu backend/app/core/config.py. KHONG SUA TAY.
\newcommand{\BangCauHinh}{%
\code{ACCESS\_TOKEN\_EXPIRE\_MINUTES} & 480 & Thời hạn access token \\
\code{REFRESH\_TOKEN\_EXPIRE\_DAYS} & 30 & Thời hạn refresh token của app mobile \\
\code{RECOGNITION\_THRESHOLD} & 0{,}45 & Ngưỡng cosine để coi là khớp danh tính \\
\code{DUPLICATE\_THRESHOLD} & 0{,}6 & Ngưỡng cảnh báo khuôn mặt trùng người khác khi đăng ký \\
\code{LATE\_AFTER\_MINUTES} & 15 & Sau bao nhiêu phút thì tính là đi muộn \\
\code{EARLY\_OPEN\_MINUTES} & 15 & Được mở phiên sớm trước giờ vào tiết \\
\code{SESSION\_GRACE\_MINUTES} & 30 & Sau giờ kết thúc bao lâu thì ngừng nhận và tự động đóng phiên \\
\code{EVIDENCE\_RETENTION\_DAYS} & 90 & Số ngày lưu ảnh bằng chứng trước khi xoá tự động \\
\code{MEDIA\_URL\_EXPIRE\_MINUTES} & 60 & Thời hạn của đường dẫn ảnh được ký điện tử (HMAC) \\
\code{ANTISPOOF\_ENABLED} & Bật & Bật hoặc tắt kiểm tra chống giả mạo \\
\code{ANTISPOOF\_THRESHOLD} & 0{,}5 & Điểm “mặt thật” tối thiểu của MiniFASNetV2 \\
\code{CONFIRM\_FRAMES} & 2 & Số khung hình phải nhận diện đúng trước khi ghi điểm danh \\
\code{SELF\_CHECKIN\_REQUIRE\_LOCATION} & Bật & Bắt buộc đối chiếu GPS khi sinh viên tự điểm danh \\
\code{SELF\_CHECKIN\_MAX\_ACCURACY\_M} & 100 & Sai số GPS tối đa chấp nhận, tính bằng mét \\
\code{ROOM\_MIN\_RADIUS\_M} & 50 & Bán kính tối thiểu khi đối chiếu GPS theo toạ độ riêng của phòng, bù sai số GPS trong nhà \\
\code{CHECKIN\_QR\_ROTATE\_SECONDS} & 30 & Chu kỳ đổi mã QR điểm danh chiếu trong phòng, tính bằng giây \\
\code{SELF\_CHECKIN\_REQUIRE\_QR\_DEFAULT} & Tắt & Mặc định phiên mới có bắt buộc quét mã QR khi tự điểm danh \\
\code{APP\_PUBLIC\_URL} &  & Địa chỉ công khai của ứng dụng, dùng cho đường link trong email và mã QR điểm danh \\
\code{NOTIFY\_EMAIL\_INTERVAL\_SECONDS} & 15 & Chu kỳ quét hàng đợi gửi email thông báo cho phụ huynh \\
\code{NOTIFY\_EMAIL\_MAX\_ATTEMPTS} & 3 & Số lần thử gửi một email thông báo trước khi đánh dấu thất bại \\
\code{ABSENCE\_WARN\_RATIO} & 0{,}2 & Tỉ lệ vắng không phép chạm ngưỡng cảnh báo cấm thi \\
\code{APP\_UTC\_OFFSET\_MINUTES} & 420 & Múi giờ gắn vào mốc thời gian trả về cho client \\
\code{LOGIN\_RATE\_LIMIT} & 30/minute & Giới hạn số lần đăng nhập theo địa chỉ IP \\
\code{LOGIN\_MAX\_FAILED\_PER\_ACCOUNT} & 5 & Số lần sai mật khẩu tối đa cho một tài khoản \\
\code{LOGIN\_FAILED\_WINDOW\_MINUTES} & 5 & Cửa sổ thời gian đếm số lần đăng nhập sai \\
\code{BACKGROUND\_JOBS\_ENABLED} & Bật & Bật hoặc tắt vòng lặp nền dọn ảnh và tự động đóng phiên \\
\code{OTP\_EXPIRE\_MINUTES} & 10 & Thời hạn mã OTP đặt lại mật khẩu \\
}
       % \BangCauHinh
% Sinh tu backend/tests/. KHONG SUA TAY.
\newcommand{\BangTest}{%
\code{test\_checkin\_qr.py} & 20 & Mã QR xoay vòng dạng link, mã 6 số dự phòng, chặn mã hết hạn \\
\code{test\_flow\_fixes.py} & 20 & Các lỗi luồng nghiệp vụ đã phát hiện và sửa \\
\code{test\_unit\_roles.py} & 20 & Vai trò theo đơn vị, phạm vi Khoa, Văn phòng Khoa và Lãnh đạo Khoa \\
\code{test\_parent\_notifications.py} & 17 & Thông báo điểm danh cho phụ huynh và hàng đợi email \\
\code{test\_admin\_crud.py} & 16 & Thêm, sửa, xoá các danh mục quản trị \\
\code{test\_rooms.py} & 16 & Danh mục phòng học, chặn trùng phòng, đối chiếu GPS theo phòng \\
\code{test\_manager\_reset\_password.py} & 14 & Cán bộ đặt lại mật khẩu: quyền từng vai trò, thu hồi phiên, huỷ OTP, mở khoá \\
\code{test\_location\_checkin.py} & 13 & Tự điểm danh kèm đối chiếu toạ độ GPS \\
\code{test\_audit\_log.py} & 12 & Ghi nhật ký tự động, che trường nhạy cảm, phạm vi xem của Lãnh đạo Khoa \\
\code{test\_consent.py} & 12 & Bắt buộc đồng ý sinh trắc: phiên bản, bản giấy, rút đồng ý, chặn nhận diện \\
\code{test\_open\_session\_rules.py} & 12 & Điều kiện mở phiên theo khung giờ của tiết học \\
\code{test\_password\_reset.py} & 12 & Quên và đặt lại mật khẩu bằng OTP \\
\code{test\_tokens\_and\_media.py} & 12 & Thu hồi token, xoay vòng refresh token, đường dẫn ảnh ký điện tử \\
\code{test\_face\_approval.py} & 11 & Duyệt và từ chối dữ liệu khuôn mặt, xoá ảnh tạm \\
\code{test\_session\_detail\_and\_records.py} & 11 & Chi tiết phiên, sửa điểm danh và nhật ký sửa \\
\code{test\_change\_password.py} & 8 & Đổi mật khẩu, thu hồi phiên đăng nhập cũ \\
\code{test\_reports\_export\_and\_warnings.py} & 8 & Cảnh báo vắng và xuất tệp Excel \\
\code{test\_antispoof.py} & 6 & Chống giả mạo, kể cả khi thiếu tệp mô hình \\
\code{test\_authorization.py} & 6 & Quyền theo vai trò trên các endpoint chính \\
\code{test\_my\_attendance\_detail.py} & 6 & Chuyên cần của sinh viên và của phụ huynh \\
\code{test\_admin\_management.py} & 5 & Quản lý tài khoản cán bộ và đơn vị \\
\code{test\_week\_schedule.py} & 5 & Thời khoá biểu theo tuần \\
\code{test\_auth.py} & 4 & Đăng nhập, sai mật khẩu, giới hạn số lần thử \\
\code{test\_confirm\_frames.py} & 4 & Xác nhận đa khung hình trước khi ghi điểm danh \\
\code{test\_reports.py} & 3 & Báo cáo chuyên cần theo lớp \\
\code{test\_session\_lifecycle.py} & 3 & Vòng đời phiên: mở, điểm danh, đóng, sinh bản ghi vắng \\
\code{test\_evidence\_cleanup.py} & 2 & Dọn ảnh bằng chứng quá hạn lưu trữ \\
}
           % \BangTest

% Ba bảng số đo thật. Mặc định là các ô chờ; chạy công cụ đo xong thì tệp
% tương ứng trong auto/ sẽ ghi đè bằng \renewcommand.
\newcommand{\BangHieuNang}{%
Phát hiện và trích đặc trưng một khuôn mặt & \cando{ms} & \cando{ms} \\
Khung hình có năm khuôn mặt & \cando{ms} & \cando{ms} \\
Truy vấn so khớp trong cơ sở dữ liệu & \cando{ms} & \cando{ms} \\
Kiểm tra chống giả mạo một khuôn mặt & \cando{ms} & \cando{ms} \\
Thời gian đáp ứng trọn vẹn một yêu cầu nhận diện & \cando{ms} & \cando{ms} \\
}
% Con so dung trong van xuoi Chuong 5. backend/tools/do_hieu_nang.py --ghi ghi de bang
% \renewcommand trong auto/bang-hieu-nang.tex; chua do thi hien o cho.
\newcommand{\HnMotMat}{\cando{ms}}
\newcommand{\HnNhieuMat}{\cando{ms}}
\newcommand{\HnSoMat}{\cando{n}}
\newcommand{\HnMoiMatThem}{\cando{ms}}
\newcommand{\HnTruyVan}{\cando{ms}}
\newcommand{\HnChongGiaMao}{\cando{ms}}
\newcommand{\HnTronVen}{\cando{ms}}
\newcommand{\HnSoVector}{\cando{n}}
\newcommand{\HnGonMotMat}{\cando{ms}}
\newcommand{\HnGonNhieuMat}{\cando{ms}}
\newcommand{\HnGiamMotMat}{\cando{\%}}
\newcommand{\HnGiamNhieuMat}{\cando{\%}}
\newcommand{\HnSoLan}{\cando{n}}
\newcommand{\HnNgayDo}{\cando{ngày đo}}
\IfFileExists{auto/bang-hieu-nang.tex}{% File nay do backend/tools/do_hieu_nang.py sinh ra. KHONG SUA TAY.
% Do luc 12/11/2026 09:16, lap 30 lan moi phep do.
\renewcommand{\BangHieuNang}{%
Phát hiện và trích đặc trưng một khuôn mặt & 380{,}9~ms & 409{,}2~ms \\
Khung hình có 6 khuôn mặt & 1438{,}0~ms & 1866{,}1~ms \\
Truy vấn so khớp trong cơ sở dữ liệu & 2{,}2~ms & 2{,}9~ms \\
Kiểm tra chống giả mạo một khuôn mặt & 3{,}0~ms & 4{,}9~ms \\
Thời gian đáp ứng trọn vẹn một yêu cầu nhận diện & 413{,}4~ms & 490{,}0~ms \\
}
\renewcommand{\HnMotMat}{381}
\renewcommand{\HnNhieuMat}{1438}
\renewcommand{\HnSoMat}{6}
\renewcommand{\HnMoiMatThem}{211}
\renewcommand{\HnTruyVan}{2{,}2}
\renewcommand{\HnChongGiaMao}{3{,}0}
\renewcommand{\HnTronVen}{413}
\renewcommand{\HnSoVector}{16}
\renewcommand{\HnGonMotMat}{278}
\renewcommand{\HnGonNhieuMat}{774}
\renewcommand{\HnGiamMotMat}{27}
\renewcommand{\HnGiamNhieuMat}{46}
\renewcommand{\HnSoLan}{30}
\renewcommand{\HnNgayDo}{12/11/2026}
}{}

% =====================================================================
%  THÔNG TIN ĐỒ ÁN
%  >>> SINH VIÊN ĐIỀN: họ tên, MSSV, lớp. Các dòng còn lại đã đúng.
% =====================================================================
\crname{BÁO CÁO ĐỒ ÁN TRÍ TUỆ NHÂN TẠO}
\cmonhoc{Học phần AIP202 \textperiodcentered{} 2 tín chỉ}
\ctname{HỆ THỐNG QUẢN LÝ SINH VIÊN TÍCH HỢP\\ĐIỂM DANH BẰNG NHẬN DIỆN KHUÔN MẶT}
\cstuname{Nguyễn Đăng Sen}
\cmssv{2351220206}
\clopsv{CNTT 23A}                % <<< KIỂM TRA LẠI: đang đoán theo khoá 23
\csCouncil{Công nghệ thông tin}
\csSupervise{TS. Đỗ Phúc Hảo}
\cttime{2026--2027}
\csdeptname{KHOA CÔNG NGHỆ THÔNG TIN}
\csplace{Đà Nẵng}

\begin{document}

\coverpage

% ---------------------------------------------------------------------
%  PHẦN ĐẦU
% ---------------------------------------------------------------------
\frontmatter
\dauPageStyle

% ---------------------------------------------------------------------
%  LỜI CAM ĐOAN
% ---------------------------------------------------------------------
\begin{declaration}

Tôi xin cam đoan đồ án \emph{\dauTenDeTai} là công trình do chính tôi thực
hiện dưới sự hướng dẫn của \dauGVHD. Các nội dung nghiên cứu, số liệu và kết
quả trình bày trong đồ án là trung thực, được lấy từ chính kho mã nguồn và
các lần chạy thử của hệ thống, và chưa từng được công bố trong bất kỳ công
trình nào khác.

Hệ thống sử dụng các mô hình và thư viện mã nguồn mở của tác giả khác, gồm
bộ mô hình InsightFace \code{buffalo\_l} (SCRFD và ArcFace)
\cite{Guo2021SCRFD, Deng2019ArcFace}, mô hình chống giả mạo MiniFASNetV2
\cite{Minivision2020Silent}, cùng các thư viện FastAPI, SQLAlchemy, pgvector
và Flutter. Toàn bộ những thành phần này đều có giấy phép cho phép sử dụng
và đã được trích dẫn nguồn gốc rõ ràng trong danh mục tài liệu tham khảo.
Phần đóng góp của tôi là toàn bộ thiết kế nghiệp vụ, cơ sở dữ liệu, quy
trình điểm danh, cơ chế phân quyền, ứng dụng đa nền tảng và bộ kiểm thử.

Trong quá trình thực hiện, tôi có sử dụng công cụ trí tuệ nhân tạo để hỗ trợ
soát lỗi và soạn thảo tài liệu. Mọi quyết định thiết kế, phần mã nguồn cốt
lõi và toàn bộ nội dung báo cáo đều do tôi kiểm chứng và chịu trách nhiệm.

Nếu phát hiện có bất kỳ sự gian lận nào, tôi xin chịu hoàn toàn trách nhiệm
trước Khoa và Nhà trường.

\vspace{1.4em}
\hfill\okyten{\dauSinhVien}{Sinh viên thực hiện}

\end{declaration}

% ---------------------------------------------------------------------
%  LỜI CẢM ƠN
% ---------------------------------------------------------------------
\begin{acknowledgments}

Em xin gửi lời cảm ơn chân thành tới \dauGVHD, người đã hướng dẫn em trong
suốt học phần \dauMonHoc. Thầy đã giúp em thu hẹp một ý tưởng ban đầu quá
rộng thành một hệ thống có phạm vi vừa sức nhưng vẫn chạy được thật, và
nhiều lần nhắc em rằng một tính năng hoàn chỉnh có giá trị hơn ba tính năng
dang dở.

Em xin cảm ơn quý thầy cô Khoa Công nghệ Thông tin, Trường Đại học Kiến trúc
Đà Nẵng đã trang bị cho em nền tảng về cơ sở dữ liệu, lập trình web và trí
tuệ nhân tạo. Những kiến thức đó là thứ em dùng trực tiếp trong từng phần
của đồ án này.

Em cũng cảm ơn các bạn trong lớp \dauLop{} đã tình nguyện cho em chụp ảnh
khuôn mặt để thử nghiệm quy trình đăng ký và nhận diện, đồng thời đóng vai
người dùng thử trong các buổi chạy thử quy trình điểm danh.

Đồ án chắc chắn còn thiếu sót. Em rất mong nhận được nhận xét của quý thầy
cô để hoàn thiện hệ thống.

\vspace{1.4em}
\hfill\okyten{\dauSinhVien}{Sinh viên thực hiện}

\end{acknowledgments}

% ---------------------------------------------------------------------
%  TÓM TẮT
% ---------------------------------------------------------------------
\begin{abstract}

Điểm danh thủ công ở giảng đường đại học vừa tốn thời gian của mỗi tiết
học, vừa không ngăn được tình trạng điểm danh hộ, và số liệu chuyên cần
thường chỉ tổng hợp được vào cuối kỳ, khi đã quá muộn để nhắc nhở sinh viên
sắp chạm ngưỡng cấm thi. Đồ án xây dựng FAMS \en{(Face Attendance Management
System)}, một hệ thống quản lý sinh viên tích hợp điểm danh bằng nhận diện
khuôn mặt, dùng được cho cả nhà trường lẫn sinh viên và phụ huynh.

Hệ thống nhận diện khuôn mặt theo hướng so khớp đặc trưng: bộ dò SCRFD tìm
khuôn mặt trong khung hình, mô hình ArcFace chuyển mỗi khuôn mặt thành vector
\FchieuVector{} chiều đã chuẩn hoá, và danh tính được xác định bằng khoảng
cách cosine tới các vector đã đăng ký, có chỉ mục HNSW của pgvector hỗ trợ
tìm kiếm. Cách này cho phép thêm một sinh viên mới chỉ bằng một ảnh, không
phải huấn luyện lại mô hình. Để chống gian lận, hệ thống kiểm tra ảnh in và
ảnh chụp lại màn hình bằng MiniFASNetV2, yêu cầu nhận diện đúng
\FsoKhungXacNhan{} khung hình mới ghi nhận, và đối chiếu toạ độ GPS khi sinh
viên tự điểm danh bằng ảnh tự chụp.

Sản phẩm gồm một dịch vụ FastAPI với \FsoEndpoint{} endpoint trên
\FsoBang{} bảng dữ liệu PostgreSQL, một ứng dụng Flutter dùng chung cho Web,
Android và iOS với \FsoManHinh{} màn hình, phân quyền cho \FsoVaiTro{} vai
trò từ quản trị viên, phòng ban, giảng viên đến sinh viên và phụ huynh. Bộ
kiểm thử tự động gồm \FsoTest{} ca chạy trên PostgreSQL thật. Kết quả cho
thấy một hệ thống điểm danh bằng khuôn mặt có thể chạy trên máy chủ không
GPU, đồng thời khẳng định rằng phần khó nhất của bài toán không nằm ở mô
hình nhận diện, mà ở các lớp phòng vệ quanh nó.

\vspace{0.8em}
\noindent\textbf{Từ khoá:} nhận diện khuôn mặt, ArcFace, chống giả mạo
khuôn mặt, tìm kiếm vector, điểm danh tự động.

\end{abstract}


\tableofcontents
\listoffigures
\listoftables
\listofalgorithms

% ---------------------------------------------------------------------
%  THÂN BÁO CÁO
% ---------------------------------------------------------------------
\mainmatter
\dauPageStyle

% ---------------------------------------------------------------------
%  CHƯƠNG 1 — MỞ ĐẦU
% ---------------------------------------------------------------------
\chapter{Mở đầu}
\label{ch:modau}

\begin{muctieu}
Chương này nêu bối cảnh của bài toán điểm danh ở giảng đường, lý do chọn
hướng giải quyết bằng nhận diện khuôn mặt, các mục tiêu cụ thể mà đồ án cam
kết đạt được, phạm vi thực hiện và bố cục của báo cáo.
\end{muctieu}

\section{Lý do chọn đề tài}

Một tiết học ở giảng đường thường bắt đầu bằng việc gọi tên hoặc chuyền tay
tờ danh sách ký. Với lớp học phần vài chục sinh viên, việc này lấy mất
khoảng năm tới mười phút đầu giờ, tức gần một phần mười thời lượng tiết học.
Tính trên toàn bộ một học kỳ, phần thời gian đó là đáng kể.

Vấn đề thứ hai nghiêm trọng hơn thời gian: điểm danh thủ công không xác thực
được danh tính. Một sinh viên có thể ký thay bạn, hoặc dạ thay khi giảng viên
gọi tên. Bản thân tờ danh sách ký cũng không để lại bằng chứng nào kiểm chứng
được về sau, nên khi có khiếu nại về kết quả chuyên cần, nhà trường gần như
không có căn cứ để đối chiếu.

Vấn đề thứ ba nằm ở khâu tổng hợp. Dữ liệu điểm danh nằm rải rác trên giấy
hoặc trong các tệp bảng tính riêng của từng giảng viên, nên thông thường chỉ
đến cuối kỳ nhà trường mới biết sinh viên nào vắng quá tỉ lệ cho phép. Khi
đó sinh viên đã mất quyền dự thi và mọi nhắc nhở đều là quá muộn. Phụ huynh
lại càng không có kênh nào để theo dõi tình hình học tập của con trong kỳ.

Nhận diện khuôn mặt là hướng giải quyết hợp lý cho ba vấn đề trên, vì khuôn
mặt là đặc điểm sinh trắc luôn mang theo người, không cần phát thêm thẻ,
không cần chạm vào thiết bị, và một khung hình chụp cả lớp có thể nhận ra
nhiều người cùng lúc. Quan trọng hơn, các mô hình nhận diện hiện đại \cite{Wang2021Survey} cho
phép thêm một người mới chỉ bằng một tấm ảnh, không phải huấn luyện lại, nên
hệ thống mở rộng được theo từng kỳ tuyển sinh. Đây cũng là bài toán cho phép
em vận dụng trọn vẹn các kiến thức về học sâu, thị giác máy tính và tìm kiếm
trên không gian vector đã học trong học phần.

\begin{luuy}
Đổi lại, dữ liệu khuôn mặt là dữ liệu sinh trắc, thuộc nhóm dữ liệu cá nhân
nhạy cảm \cite{Luat912025, ND3562025}. Một hệ thống điểm danh bằng khuôn mặt chỉ có ý nghĩa thực tiễn nếu
đi kèm cơ chế xin phép, phân quyền truy cập và quyền xoá dữ liệu. Ràng buộc
này được coi là yêu cầu bắt buộc của đồ án, chứ không phải phần làm thêm.
\end{luuy}

\section{Mục tiêu của đồ án}
\label{sec:muctieu}

Đồ án đặt ra năm mục tiêu cụ thể, mỗi mục tiêu đều kiểm chứng được và sẽ
được đối chiếu lại trong Chương~\ref{ch:ketqua}:

\begin{enumerate}[label=\textbf{MT\arabic*.}, leftmargin=3.0em, itemsep=3pt]
  \item \textbf{Xây dựng được quy trình nhận diện khuôn mặt hoàn chỉnh} từ
        đăng ký tới điểm danh, chạy được trên máy chủ chỉ có CPU, thời gian
        xử lý một khung hình đủ nhanh để dùng trực tiếp trong lớp.
  \item \textbf{Chống được ba kiểu gian lận thường gặp}: điểm danh bằng ảnh
        in hoặc ảnh chụp lại màn hình, nhận nhầm danh tính do một khung hình
        xấu, và điểm danh từ xa khi không có mặt tại lớp.
  \item \textbf{Số hoá trọn vẹn nghiệp vụ quản lý chuyên cần}: hồ sơ sinh
        viên, lớp học phần, lịch học, phiên điểm danh, sửa điểm danh có lưu
        vết, cảnh báo vắng, thông báo cho phụ huynh và xuất báo cáo, kèm nhật
        ký hệ thống ghi lại mọi thao tác quản lý.
  \item \textbf{Phân quyền đúng mô hình tổ chức của trường} cho quản trị
        viên, Phòng Đào tạo, Phòng Công tác sinh viên, văn phòng Khoa, lãnh
        đạo Khoa, ban ngành, giảng viên, sinh viên và phụ huynh, trong đó cán
        bộ Khoa chỉ thấy dữ liệu thuộc Khoa mình và người nhập liệu tách khỏi
        người phê duyệt.
  \item \textbf{Cung cấp một ứng dụng dùng chung cho Web, Android và iOS}
        từ một mã nguồn duy nhất, kèm bộ kiểm thử tự động chạy trên cơ sở dữ
        liệu thật để bảo vệ các luồng nghiệp vụ quan trọng.
\end{enumerate}

\section{Phạm vi và đối tượng}

\textbf{Về bài toán trí tuệ nhân tạo}, đồ án sử dụng mô hình nhận diện khuôn
mặt đã được tiền huấn luyện thay vì huấn luyện một mô hình mới từ đầu. Lý do
được phân tích ở Mục~\ref{sec:vosaokhonghuanluyen}. Phần đóng góp của đồ án
nằm ở quy trình bao quanh mô hình: kiểm tra chất lượng ảnh, chống giả mạo,
xác nhận đa khung hình, so khớp theo phạm vi lớp và đối chiếu vị trí.

\textbf{Về nghiệp vụ}, đồ án số hoá quy trình điểm danh và quản lý chuyên
cần. Hệ thống không quản lý điểm học phần, học phí, đăng ký tín chỉ hay
nhân sự, và không kết nối với phần mềm quản lý đào tạo sẵn có của trường.

\textbf{Về người dùng}, hệ thống phục vụ \FsoVaiTro{} vai trò được liệt kê ở
Mục~\ref{sec:phanquyen}. Ứng dụng web hướng tới cán bộ và giảng viên, còn
ứng dụng di động hướng tới sinh viên và phụ huynh.

\textbf{Những gì đồ án không làm}: không nhận diện khi sinh viên đeo khẩu
trang che quá nửa khuôn mặt; không xử lý được trường hợp anh chị em sinh đôi
cùng trứng bằng riêng khuôn mặt; không triển khai trên hạ tầng thật của nhà
trường; và chưa kiểm định mô hình trên tập dữ liệu khuôn mặt người Việt quy
mô lớn.

\section{Phương pháp thực hiện}

Đồ án được thực hiện theo hướng phát triển lặp \cite{Larman2003Iterative}. Mỗi vòng lặp gồm bốn bước:
chọn một luồng nghiệp vụ, cài đặt luồng đó ở cả phía dịch vụ lẫn ứng dụng,
viết kiểm thử tự động cho luồng, rồi chạy thử với dữ liệu mẫu để tìm lỗi
thực tế. Cách làm này giữ cho hệ thống luôn ở trạng thái chạy được, thay vì
dồn toàn bộ phần tích hợp vào cuối kỳ.

Toàn bộ mã nguồn được quản lý bằng \code{git}. Tại thời điểm viết báo cáo,
kho mã nguồn có \FsoCommit{} commit, bản được mô tả trong báo cáo này là
commit \code{\FcommitCuoi} ngày \FngayCommit.

\begin{ghinho}
Mọi con số về quy mô hệ thống trong báo cáo này, từ số bảng dữ liệu tới số
ca kiểm thử, đều được sinh tự động từ chính kho mã nguồn bằng công cụ mô tả
ở mục C của phần Phụ lục, nên không có số nào là ước lượng bằng tay.
\end{ghinho}

\section{Bố cục báo cáo}

Phần còn lại của báo cáo được tổ chức như sau. Chương~\ref{ch:coso} trình
bày cơ sở lý thuyết về nhận diện khuôn mặt, chống giả mạo, tìm kiếm trên
không gian vector và các công nghệ được sử dụng. Chương~\ref{ch:thietke}
phân tích yêu cầu và thiết kế kiến trúc, cơ sở dữ liệu, giao diện lập trình
cùng mô hình phân quyền. Chương~\ref{ch:caidat} mô tả quá trình cài đặt các
thành phần cốt lõi và giao diện người dùng. Chương~\ref{ch:ketqua} trình bày
kết quả đạt được, kết quả kiểm thử, đánh giá so với mục tiêu ban đầu, rồi
tổng kết, nêu bài học rút ra và hướng phát triển.

% ---------------------------------------------------------------------
%  CHƯƠNG 2 — CƠ SỞ LÝ THUYẾT VÀ CÔNG NGHỆ
% ---------------------------------------------------------------------
\chapter{Cơ sở lý thuyết và công nghệ}
\label{ch:coso}

\begin{muctieu}
Chương này trình bày nền tảng lý thuyết mà hệ thống dựa vào: cách một khuôn
mặt được phát hiện và chuyển thành vector đặc trưng, cách so khớp danh tính
bằng khoảng cách cosine, nguyên lý chống giả mạo khuôn mặt, kỹ thuật tìm
kiếm trên không gian vector, cùng các công nghệ được chọn để hiện thực hoá.
\end{muctieu}

\section{Bài toán nhận diện khuôn mặt}

Cần phân biệt ba bài toán thường bị gộp làm một \cite{Jain2004Biometric, Wang2021Survey}:

\begin{dinhnghia}[Ba bài toán trên khuôn mặt]
\textbf{Phát hiện khuôn mặt} \en{(face detection)} trả lời câu hỏi trong ảnh
có khuôn mặt nào và nằm ở đâu. \textbf{Xác minh khuôn mặt} \en{(face
verification)} trả lời hai khuôn mặt có phải cùng một người hay không.
\textbf{Nhận diện khuôn mặt} \en{(face recognition hay identification)} trả
lời khuôn mặt này là ai trong số những người đã biết.
\end{dinhnghia}

Điểm danh trong lớp là bài toán nhận diện trên tập đóng có giới hạn: hệ thống
chỉ cần xác định khuôn mặt trong khung hình thuộc về sinh viên nào trong danh
sách lớp, chứ không phải toàn bộ cơ sở dữ liệu của trường. Việc thu hẹp
không gian tìm kiếm theo lớp học phần vừa làm tăng tốc độ, vừa giảm mạnh xác
suất nhận nhầm, và được hệ thống khai thác triệt để ở Mục~\ref{sec:sokhop}.

Quy trình nhận diện hiện đại gồm bốn bước nối tiếp nhau \cite{Wang2021Survey} như Hình~\ref{fig:pipeline}.

\begin{figure}[H]
\centering
\begin{tikzpicture}[node distance=1.15cm and 0.85cm, every node/.append style={align=center}]
  \node (anh)   [io]      {Khung hình\\ từ camera};
  \node (detect)[process, right=of anh]    {Phát hiện\\ SCRFD};
  \node (align) [process, right=of detect] {Căn chỉnh\\ theo 5 điểm mốc};
  \node (embed) [process, right=of align]  {Trích đặc trưng\\ ArcFace};
  \node (match) [decision, below=of align] {So khớp cosine};
  \node (out)   [startstop, right=of match, xshift=0.6cm] {Danh tính\\ và độ tin cậy};
  \draw [flowarr] (anh)    -- (detect);
  \draw [flowarr] (detect) -- (align);
  \draw [flowarr] (align)  -- (embed);
  \draw [flowarr] (embed)  |- (match);
  \draw [flowarr] (match)  -- (out);
\end{tikzpicture}
\caption{Bốn bước của quy trình nhận diện khuôn mặt trong FAMS.}
\label{fig:pipeline}
\end{figure}

\section{Phát hiện khuôn mặt bằng SCRFD}

Bước đầu tiên là tìm vị trí mọi khuôn mặt trong khung hình. Các phương pháp
cổ điển dựa trên đặc trưng Haar và bộ phân loại tầng \cite{Viola2001Rapid}
chạy nhanh nhưng suy giảm rõ rệt khi khuôn mặt nghiêng hoặc thiếu sáng, tức
đúng những điều kiện thường gặp trong giảng đường.

Hệ thống dùng SCRFD \cite{Guo2021SCRFD}, một bộ dò khuôn mặt một giai đoạn
được thiết kế để phân bổ chi phí tính toán hợp lý giữa các tầng đặc trưng,
nhờ đó đạt độ chính xác cao trên khuôn mặt nhỏ mà vẫn nhẹ. Với ảnh chụp cả
lớp, khuôn mặt của sinh viên ngồi cuối phòng chiếm rất ít điểm ảnh, nên khả
năng phát hiện khuôn mặt nhỏ là tiêu chí quyết định lựa chọn này. Ngoài khung
bao, SCRFD trả về năm điểm mốc gồm hai mắt, chóp mũi và hai khoé miệng, dùng
cho bước căn chỉnh tiếp theo.

Căn chỉnh \en{(alignment)} là phép biến đổi affine đưa năm điểm mốc về vị trí
chuẩn, cho ra một ảnh khuôn mặt kích thước cố định \cite{Wang2021Survey}. Bước này quan trọng vì mô
hình trích đặc trưng được huấn luyện trên các khuôn mặt đã căn chỉnh; bỏ qua
nó thì cùng một người ở hai góc chụp khác nhau sẽ cho hai vector lệch nhau
đáng kể.

\section{Biểu diễn khuôn mặt bằng ArcFace}
\label{sec:arcface}

Thay vì phân loại trực tiếp ảnh khuôn mặt thành các lớp là từng người, hướng
tiếp cận hiện đại học một \term{hàm nhúng} \en{(embedding)} ánh xạ ảnh khuôn
mặt vào không gian vector, sao cho hai ảnh của cùng một người nằm gần nhau và
hai người khác nhau nằm xa nhau. Ý tưởng này được phổ biến bởi FaceNet
\cite{Schroff2015FaceNet} và được ArcFace \cite{Deng2019ArcFace} cải tiến.

ArcFace bổ sung một biên góc cộng tính vào hàm mất mát softmax. Gọi
$\theta_{y_i}$ là góc giữa vector đặc trưng của mẫu thứ $i$ và vector trọng
số của lớp đúng $y_i$, $s$ là bán kính cầu và $m$ là biên góc, hàm mất mát
có dạng \cite{Deng2019ArcFace}:

\begin{equation}
\label{eq:arcface}
L = -\frac{1}{N}\sum_{i=1}^{N}
\log \frac{e^{s\cos(\theta_{y_i}+m)}}
{e^{s\cos(\theta_{y_i}+m)} + \sum_{j \neq y_i} e^{s\cos\theta_j}}
\end{equation}

Việc cộng thêm $m$ vào góc của lớp đúng buộc mô hình phải kéo mẫu lại gần
tâm lớp của nó một khoảng chặt hơn mức tối thiểu cần để phân loại đúng. Hệ
quả là các cụm đặc trưng của từng người tách bạch hơn trên mặt cầu đơn vị, và
đó chính là tính chất mà bài toán nhận diện cần.

Mô hình được dùng là bộ \code{buffalo\_l} của thư viện InsightFace
\cite{InsightFace2023}, sinh ra vector \FchieuVector{} chiều đã chuẩn hoá về
độ dài đơn vị.

\begin{ghinho}
Vector đặc trưng đã chuẩn hoá nên tích vô hướng của hai vector chính bằng
cosine của góc giữa chúng. Nhờ vậy phép so khớp danh tính chỉ còn là một
phép nhân vô hướng, rất rẻ về mặt tính toán.
\end{ghinho}

\subsection{So khớp bằng độ tương tự cosine}
\label{sec:sokhop}

Với hai vector đặc trưng $\mathbf{u}$ và $\mathbf{v}$, độ tương tự cosine và
khoảng cách cosine được định nghĩa \cite{Manning2008IR}:

\begin{equation}
\label{eq:cosine}
\mathrm{sim}(\mathbf{u}, \mathbf{v}) =
\frac{\mathbf{u} \cdot \mathbf{v}}{\lVert \mathbf{u}\rVert \, \lVert \mathbf{v}\rVert},
\qquad
d(\mathbf{u}, \mathbf{v}) = 1 - \mathrm{sim}(\mathbf{u}, \mathbf{v})
\end{equation}

Giá trị $\mathrm{sim}$ nằm trong đoạn $[-1, 1]$, càng gần $1$ thì hai khuôn
mặt càng giống nhau. Hệ thống coi là cùng một người khi độ tương tự vượt
ngưỡng $\tau = \FnguongNhanDien$. Ngưỡng này điều tiết trực tiếp cán cân giữa
hai loại lỗi \cite{Jain2004Biometric, ISO19795}:

\begin{itemize}[leftmargin=1.6em, itemsep=2pt]
  \item \textbf{Chấp nhận sai} \en{(false accept)}: nhận nhầm sinh viên này
        thành sinh viên khác, dẫn tới điểm danh cho người vắng mặt. Đây là
        lỗi nguy hiểm vì phá hỏng ý nghĩa của cả hệ thống.
  \item \textbf{Từ chối sai} \en{(false reject)}: không nhận ra sinh viên
        đang có mặt. Lỗi này gây phiền nhưng còn khắc phục được, vì giảng
        viên có thể xác nhận tay.
\end{itemize}

Vì hai loại lỗi có hậu quả không cân xứng, hệ thống ưu tiên hạ thấp tỉ lệ
chấp nhận sai, chấp nhận đánh đổi bằng một mức từ chối sai cao hơn.

\subsection{Vì sao không tự huấn luyện mô hình}
\label{sec:vosaokhonghuanluyen}

Một câu hỏi hợp lý là tại sao đồ án không huấn luyện một mạng phân loại trên
chính ảnh sinh viên của trường, mà dùng lại mô hình đã huấn luyện sẵn theo
hướng học chuyển giao \en{(transfer learning)} \cite{Pan2010Transfer}. Có ba lý do:

\begin{enumerate}[leftmargin=1.8em, itemsep=2pt]
  \item \textbf{Tập người dùng luôn thay đổi.} Mỗi kỳ có sinh viên mới nhập
        học và sinh viên tốt nghiệp. Mô hình phân loại buộc phải huấn luyện
        lại mỗi lần danh sách đổi, còn hàm nhúng chỉ cần thêm một vector.
  \item \textbf{Dữ liệu cho mỗi người quá ít.} Mỗi sinh viên chỉ đăng ký
        vài ảnh, không đủ để huấn luyện một lớp phân loại ổn định, trong khi
        ArcFace được huấn luyện trên hàng triệu ảnh của hàng chục nghìn danh
        tính \cite{Deng2019ArcFace} nên đã học được đặc trưng tổng quát.
  \item \textbf{Chi phí và rủi ro.} Huấn luyện lại đòi hỏi GPU và một tập
        dữ liệu khuôn mặt sinh viên đủ lớn, tức phải thu thập thêm dữ liệu
        sinh trắc, đi ngược nguyên tắc thu thập tối thiểu \cite{GDPR2016} đã nêu ở
        Chương~\ref{ch:modau}.
\end{enumerate}

\section{Chống giả mạo khuôn mặt}
\label{sec:chonggiamao}

Một hệ thống chỉ so khớp đặc trưng sẽ chấp nhận cả ảnh in lẫn ảnh hiển thị
trên màn hình điện thoại, vì những ảnh đó cho vector đặc trưng gần như giống
hệt người thật. Đây là kiểu tấn công trình chiếu \en{(presentation attack)} \cite{ISO30107, Ramachandra2017Survey},
và là lỗ hổng rõ ràng nhất của điểm danh bằng khuôn mặt.

Phát hiện sự sống \en{(liveness detection)} chia làm hai nhóm \cite{Ramachandra2017Survey}. Nhóm chủ động
yêu cầu người dùng thực hiện hành vi như chớp mắt hay quay đầu; nhóm bị động
chỉ phân tích chính khung hình tĩnh. Hệ thống chọn hướng bị động với mô hình
MiniFASNetV2 thuộc bộ Silent-Face-Anti-Spoofing \cite{Minivision2020Silent},
vì hướng chủ động không dùng được khi giảng viên quét cả lớp bằng một khung
hình.

MiniFASNetV2 khai thác những dấu vết mà một bề mặt phẳng để lại \cite{Ramachandra2017Survey}: phản xạ ánh
sáng khác da thật, hiện tượng răng cưa của lưới điểm ảnh khi chụp lại màn
hình, viền giấy hoặc viền thiết bị, và sự thiếu vắng kết cấu vi mô của da.
Mô hình nhận ảnh khuôn mặt đã mở rộng khung bao theo hệ số $2{,}7$ để lấy
thêm bối cảnh xung quanh, đưa về kích thước $80 \times 80$ điểm ảnh, rồi cho
ra điểm tin cậy cho ba lớp, trong đó lớp \term{mặt thật} phải vượt ngưỡng
$\FnguongThat$ thì khung hình mới được chấp nhận.

\begin{luuy}
Không có mô hình chống giả mạo nào đạt độ chính xác tuyệt đối. Vì vậy hệ
thống không coi đây là lớp phòng vệ duy nhất, mà kết hợp thêm xác nhận đa
khung hình, đối chiếu vị trí và quyền xác nhận cuối cùng của giảng viên, như
trình bày ở Mục~\ref{sec:baolop}.
\end{luuy}

\section{Tìm kiếm trên không gian vector}

Khi số vector đã đăng ký tăng lên, việc so sánh tuần tự với từng vector trở
nên tốn kém. Bài toán này gọi là tìm láng giềng gần nhất \en{(nearest
neighbour search)}. Với dữ liệu nhiều chiều, các cấu trúc chỉ mục cổ điển
không còn hiệu quả, nên thực tế dùng tìm kiếm gần đúng \cite{Indyk1998ANN}, đổi một phần rất nhỏ
độ chính xác lấy tốc độ cao hơn nhiều lần.

Hệ thống dùng chỉ mục HNSW \cite{Malkov2020HNSW}, một đồ thị thế giới nhỏ
phân tầng: tầng trên thưa để nhảy nhanh tới vùng cần tìm, tầng dưới dày để
tinh chỉnh kết quả. Điểm thuận lợi về mặt kỹ thuật là chỉ mục này có sẵn
trong pgvector \cite{Pgvector2024}, một phần mở rộng của PostgreSQL cho phép
lưu vector ngay trong bảng dữ liệu và viết truy vấn so khớp bằng SQL thông
thường. Nhờ đó hệ thống không cần thêm một cơ sở dữ liệu vector riêng, và
điều kiện lọc theo lớp học phần hay theo trạng thái duyệt được đẩy xuống cho
cơ sở dữ liệu xử lý cùng lúc với phép so khớp.

\section{Công nghệ sử dụng}

\begin{table}[H]
\centering
\caption{Các công nghệ chính và lý do lựa chọn.}
\label{tab:congnghe}
\small
\begin{tabularx}{\textwidth}{@{}>{\bfseries}l >{\raggedright\arraybackslash}p{3.1cm} X@{}}
\toprule
Thành phần & Công nghệ & Lý do chọn \\
\midrule
Dịch vụ & FastAPI, SQLAlchemy 2.0 & Kiểu dữ liệu tường minh, tự sinh tài liệu
  API, hệ sinh thái Python sẵn có thư viện thị giác máy tính \\
Cơ sở dữ liệu & PostgreSQL 16, pgvector \cite{Pgvector2024} & Lưu dữ liệu quan hệ và vector đặc
  trưng trong cùng một nơi, hỗ trợ chỉ mục HNSW \\
Suy luận mô hình & ONNX Runtime (CPU) & Chạy được trên máy chủ không GPU,
  không phụ thuộc khung học sâu lúc triển khai \\
Ứng dụng & Flutter & Một mã nguồn cho Web, Android và iOS, truy cập được
  camera và định vị \\
Triển khai & Docker Compose, nginx & Dựng lại toàn bộ hệ thống bằng một
  lệnh, tách biệt môi trường \\
Kiểm thử & pytest trên PostgreSQL thật & Kiểu dữ liệu vector không mô phỏng
  được bằng cơ sở dữ liệu nhúng \\
\bottomrule
\end{tabularx}
\end{table}

{\small
\begin{xltabular}{\textwidth}{@{}>{\bfseries\raggedright\arraybackslash\leavevmode}p{2.6cm} >{\raggedright\arraybackslash\leavevmode}p{5.6cm} >{\raggedright\arraybackslash\leavevmode}X@{}}
\caption{Giải thích chi tiết Bảng~\ref{tab:congnghe}: mỗi công nghệ nằm ở đâu trong mã nguồn.}
\label{tab:congnghe-gt}\\
\toprule
Lớp & Nằm ở & Ghi chú \\
\midrule
\endfirsthead
\toprule
Lớp & Nằm ở & Ghi chú \\
\midrule
\endhead
\bottomrule
\endfoot
Dịch vụ & \code{backend/app/api}, \code{backend/app/services}, \code{backend/app/schemas.py} & FastAPI tự kiểm tra dữ liệu vào theo kiểu đã khai báo và trả lỗi 422 khi sai, nên từng endpoint không phải tự kiểm tra từng trường \\
Cơ sở dữ liệu & \code{models.py} trong \code{backend/app/db}; các tệp di trú trong \code{backend/alembic} & Lược đồ chỉ đổi qua di trú Alembic; cột vector có chỉ mục HNSW \\
Suy luận mô hình & \code{face\_service.py}, \code{antispoof\_service.py} trong \code{backend/app/services} & Mô hình nạp ở lần dùng đầu tiên rồi giữ trong bộ nhớ cho các yêu cầu sau \\
Ứng dụng & \code{mobile/lib} & Cùng một mã nguồn build ra bản Web, Android và iOS \\
Triển khai & \code{docker-compose.yml}, \code{deploy/} & Một lệnh dựng cơ sở dữ liệu, máy chủ và giao diện \\
Kiểm thử & \code{backend/tests}, \code{mobile/test} & Phía máy chủ chạy trên PostgreSQL thật; phía ứng dụng dùng \code{flutter test} \\
\end{xltabular}}

\section{Các hệ thống liên quan}

Điểm danh tự động đã có nhiều lời giải thương mại và học thuật. Bảng
\ref{tab:sosanh} đối chiếu FAMS với ba nhóm giải pháp phổ biến.

{\small
\begin{xltabular}{\textwidth}{@{}>{\raggedright\arraybackslash}p{3.0cm} X X X@{}}
\caption{So sánh FAMS với các giải pháp điểm danh khác.}
\label{tab:sosanh}\\
\toprule
Tiêu chí & Điểm danh thủ công & Thẻ từ hoặc mã QR & FAMS \\
\midrule
\endfirsthead
\toprule
Tiêu chí & Điểm danh thủ công & Thẻ từ hoặc mã QR & FAMS \\
\midrule
\endhead
\bottomrule
\endfoot
Thời gian mỗi tiết & 5--10 phút & 1--2 phút & Dưới 1 phút, quét cả lớp \\
Chống điểm danh hộ & Không & Không, vì thẻ và mã chuyển tay được & Có, dựa
  trên đặc điểm sinh trắc \\
Bằng chứng kiểm chứng & Chữ ký giấy & Nhật ký quẹt thẻ & Ảnh bằng chứng kèm
  nhật ký sửa đổi \\
Chi phí thiết bị & Không & Đầu đọc cho mỗi phòng & Máy tính hoặc điện thoại
  sẵn có \\
Rủi ro về dữ liệu & Thấp & Trung bình & Cao, vì là dữ liệu sinh trắc, nên
  cần cơ chế đồng ý và phân quyền \\
\end{xltabular}}

{\small
\begin{xltabular}{\textwidth}{@{}>{\bfseries\raggedright\arraybackslash\leavevmode}p{2.6cm} >{\raggedright\arraybackslash\leavevmode}p{4.6cm} >{\raggedright\arraybackslash\leavevmode}X@{}}
\caption{Giải thích chi tiết Bảng~\ref{tab:sosanh}: ý nghĩa từng tiêu chí và căn cứ cho cột FAMS.}
\label{tab:sosanh-gt}\\
\toprule
Tiêu chí & Nghĩa là gì & Căn cứ cho cột FAMS \\
\midrule
\endfirsthead
\toprule
Tiêu chí & Nghĩa là gì & Căn cứ cho cột FAMS \\
\midrule
\endhead
\bottomrule
\endfoot
Thời gian mỗi tiết & Thời gian giảng viên dành cho việc điểm danh trong một tiết & Ước lượng từ thời gian xử lý mỗi khung hình ở Bảng~\ref{tab:hieunang}; chưa đo trên một lớp học thật \\
Chống điểm danh hộ & Có ngăn được một người điểm danh thay người khác hay không & Khuôn mặt không chuyển tay được; các lớp kiểm tra ở Bảng~\ref{tab:phongve} \\
Bằng chứng kiểm chứng & Căn cứ để xem lại khi có khiếu nại & Ảnh bằng chứng lưu \FngayLuuAnh{} ngày và nhật ký sửa điểm danh có lý do bắt buộc \\
Chi phí thiết bị & Thiết bị phải mua thêm cho mỗi phòng & Dùng webcam hoặc điện thoại sẵn có; máy chủ chỉ cần CPU \\
Rủi ro về dữ liệu & Mức nhạy cảm của dữ liệu thu thập & Dữ liệu sinh trắc, nên đồng ý là bắt buộc (Mục~\ref{sec:dongy}) và mọi truy cập đi qua phân quyền \\
\end{xltabular}}

Nhận xét rút ra là ưu thế của hướng nhận diện khuôn mặt chỉ giữ được nếu hệ
thống xử lý đúng ba việc: chống giả mạo, phân quyền truy cập ảnh, và cho
người dùng quyền kiểm soát dữ liệu sinh trắc của mình. Ba việc này định hình
phần lớn thiết kế được trình bày ở Chương~\ref{ch:thietke}.

% ---------------------------------------------------------------------
%  CHƯƠNG 3 — PHÂN TÍCH VÀ THIẾT KẾ HỆ THỐNG
% ---------------------------------------------------------------------
\chapter{Phân tích và thiết kế hệ thống}
\label{ch:thietke}

\begin{muctieu}
Chương này phân tích yêu cầu và ca sử dụng của từng nhóm người dùng, trình
bày kiến trúc tổng thể, thiết kế cơ sở dữ liệu và giao diện lập trình, mô hình
phân quyền theo chức năng và phạm vi dữ liệu, các lớp phòng vệ chống gian lận
điểm danh, cơ chế đồng ý sinh trắc học và nhật ký hệ thống.
\end{muctieu}

\section{Phân tích yêu cầu}

\subsection{Tác nhân của hệ thống}

Hệ thống phục vụ \FsoVaiTro{} vai trò, phản ánh đúng cơ cấu tổ chức của một
trường đại học chứ không gộp tất cả cán bộ thành một vai trò quản trị duy
nhất. Riêng cấp Khoa được tách thành hai vai trò: văn phòng Khoa làm công
việc nhập liệu, còn lãnh đạo Khoa giám sát và phê duyệt.

{\small
\begin{xltabular}{\textwidth}{@{}>{\bfseries}p{3.4cm} X@{}}
\caption{Các tác nhân và nhu cầu chính.}
\label{tab:tacnhan}\\
\toprule
Tác nhân & Nhu cầu chính \\
\midrule
\endfirsthead
\toprule
Tác nhân & Nhu cầu chính \\
\midrule
\endhead
\bottomrule
\endfoot
Quản trị viên & Quản trị toàn hệ thống: đơn vị, tài khoản cán bộ, dữ liệu
  gốc, xem nhật ký hệ thống, xử lý sự cố \\
Phòng Đào tạo & Tổ chức đào tạo toàn trường: lớp học phần, lịch học, phòng
  học, ghi danh, phân công giảng dạy \\
Phòng Công tác sinh viên & Hồ sơ sinh viên và phụ huynh, ghi nhận đồng ý sinh
  trắc học, duyệt dữ liệu khuôn mặt, duyệt nghỉ có phép \\
Văn phòng Khoa & Nhập liệu của Khoa: lớp học phần, lịch học, ghi danh, phân
  công giảng viên trong Khoa mình \\
Lãnh đạo Khoa & Giám sát chuyên cần của Khoa, duyệt nghỉ có phép, xem nhật
  ký thao tác của Khoa; không nhập liệu \\
Ban ngành & Xem số liệu chuyên cần tổng hợp, không xem thông tin cá nhân \\
Giảng viên & Mở phiên, quét camera, chiếu mã QR, xác nhận tay, sửa điểm danh
  lớp mình dạy \\
Sinh viên & Đồng ý hoặc rút đồng ý sinh trắc học, đăng ký khuôn mặt, tự điểm
  danh, theo dõi chuyên cần \\
Phụ huynh & Theo dõi chuyên cần của con, nhận thông báo mỗi lần con được
  điểm danh \\
\end{xltabular}}

{\small
\begin{xltabular}{\textwidth}{@{}>{\bfseries\raggedright\arraybackslash\leavevmode}p{3.6cm} >{\raggedright\arraybackslash\leavevmode}X >{\raggedright\arraybackslash\leavevmode}p{4.4cm}@{}}
\caption{Giải thích chi tiết Bảng~\ref{tab:tacnhan}: phạm vi dữ liệu và thiết bị của từng tác nhân.}
\label{tab:tacnhan-gt}\\
\toprule
Tác nhân & Thấy dữ liệu của ai & Dùng chủ yếu trên \\
\midrule
\endfirsthead
\toprule
Tác nhân & Thấy dữ liệu của ai & Dùng chủ yếu trên \\
\midrule
\endhead
\bottomrule
\endfoot
Quản trị viên & Toàn trường & Web \\
Phòng Đào tạo & Toàn trường & Web \\
Phòng Công tác sinh viên & Toàn trường & Web \\
Văn phòng Khoa & Chỉ sinh viên, giảng viên, lớp thuộc Khoa mình & Web \\
Lãnh đạo Khoa & Chỉ Khoa mình, kể cả nhật ký thao tác liên quan tới Khoa & Web \\
Ban ngành & Chỉ số liệu tổng hợp, không có tên hay ảnh & Web \\
Giảng viên & Các lớp mình dạy & Web, máy tính trong phòng để quét camera và chiếu mã QR \\
Sinh viên & Dữ liệu của chính mình & Điện thoại \\
Phụ huynh & Con đã được Phòng Công tác sinh viên liên kết & Điện thoại \\
\end{xltabular}}

\subsection{Ca sử dụng theo từng tác nhân}
\label{sec:casudung}

Mỗi tác nhân chỉ được giao những ca sử dụng \cite{OMG2017UML} gắn với nhiệm vụ của mình trong
tổ chức. Bảng~\ref{tab:casudung} liệt kê đầy đủ; ngoài các ca này, mọi tài
khoản đều dùng chung ba ca cơ bản là đăng nhập, đổi mật khẩu và quên mật khẩu
bằng mã OTP.

{\footnotesize
\begin{xltabular}{\textwidth}{@{}>{\bfseries\raggedright\arraybackslash}p{2.7cm} >{\raggedright\arraybackslash}X@{}}
\caption{Ca sử dụng của từng tác nhân.}
\label{tab:casudung}\\
\toprule
Tác nhân & Ca sử dụng \\
\midrule
\endfirsthead
\toprule
Tác nhân & Ca sử dụng \\
\midrule
\endhead
\bottomrule
\endfoot
Quản trị viên & UC01 Quản lý đơn vị, tài khoản cán bộ, đặt lại mật khẩu mọi tài khoản;
  UC02 Quản lý mọi danh mục khi cần thay phòng ban;
  UC03 Xem nhật ký hệ thống toàn trường \\
Phòng Đào tạo & UC04 Quản lý lớp học phần, ghi danh, phân công giảng viên;
  UC05 Xếp lịch học vào phòng, chặn trùng phòng;
  UC06 Quản lý phòng học và địa điểm điểm danh;
  UC07 Quản lý hồ sơ giảng viên;
  UC08 Xem báo cáo chuyên cần toàn trường \\
Phòng CTSV & UC09 Quản lý hồ sơ sinh viên, phụ huynh, liên kết và đặt lại mật khẩu của họ;
  UC10 Ghi nhận đồng ý sinh trắc theo bản giấy, rút đồng ý theo đơn;
  UC11 Duyệt ảnh khuôn mặt;
  UC12 Duyệt nghỉ có phép toàn trường \\
Văn phòng Khoa & UC13 Quản lý lớp học phần, lịch học, ghi danh của Khoa;
  UC14 Xem sinh viên, giảng viên, điểm danh, báo cáo của Khoa \\
Lãnh đạo Khoa & UC15 Giám sát chuyên cần và cảnh báo vắng của Khoa;
  UC16 Duyệt nghỉ có phép trong Khoa;
  UC17 Xem nhật ký thao tác của Khoa \\
Ban ngành & UC18 Xem thống kê tổng hợp không định danh \\
Giảng viên & UC19 Mở và đóng phiên điểm danh;
  UC20 Quét camera nhận diện cả lớp;
  UC21 Chiếu mã QR điểm danh;
  UC22 Xác nhận tay, sửa điểm danh kèm lý do;
  UC23 Xem báo cáo lớp mình dạy \\
Sinh viên & UC24 Đồng ý hoặc rút đồng ý sinh trắc học;
  UC25 Đăng ký khuôn mặt;
  UC26 Tự điểm danh bằng khuôn mặt, vị trí và mã QR;
  UC27 Xem lịch học và chuyên cần \\
Phụ huynh & UC28 Xem chuyên cần của con;
  UC29 Nhận thông báo điểm danh trong ứng dụng và qua email;
  UC30 Chọn loại thông báo muốn nhận \\
\end{xltabular}}

{\footnotesize
\begin{xltabular}{\textwidth}{@{}>{\bfseries\raggedright\arraybackslash\leavevmode}p{2.6cm} >{\raggedright\arraybackslash\leavevmode}p{3.2cm} >{\raggedright\arraybackslash\leavevmode}X >{\raggedright\arraybackslash\leavevmode}X@{}}
\caption{Giải thích chi tiết Bảng~\ref{tab:casudung}: đặc tả một số ca sử dụng tiêu biểu.}
\label{tab:casudung-gt}\\
\toprule
Ca sử dụng & Điều kiện trước & Luồng chính & Kết quả và ngoại lệ \\
\midrule
\endfirsthead
\toprule
Ca sử dụng & Điều kiện trước & Luồng chính & Kết quả và ngoại lệ \\
\midrule
\endhead
\bottomrule
\endfoot
UC19 Mở và đóng phiên (giảng viên) & Có lịch học của lớp mình hôm nay; trong khoảng từ \FphutMoSom{} phút trước giờ vào tới giờ kết thúc & Chọn buổi ở mục lịch dạy hôm nay, bấm Mở phiên, tuỳ chọn bật yêu cầu mã QR, nhận điểm danh rồi bấm Đóng phiên & Lúc đóng, sinh viên chưa có bản ghi thành Vắng và cảnh báo được tính lại; quên đóng thì hệ thống tự đóng sau \FphutAnHan{} phút. Mở ngoài khung giờ bị từ chối \\
UC20 Quét camera nhận diện cả lớp (giảng viên) & Phiên đang mở & Màn quét chụp định kỳ, mặc định 2 giây một lần, gửi khung hình lên; máy chủ chỉ so khớp với sinh viên của lớp & Sinh viên được nhận ra đủ \FsoKhungXacNhan{} lần thì được ghi Có mặt, hoặc Muộn nếu quá \FphutMuon{} phút; khuôn mặt nghi giả mạo được đánh dấu và không được ghi \\
UC26 Tự điểm danh (sinh viên) & Đã ghi danh, đã đồng ý sinh trắc học, có khuôn mặt đã duyệt, phiên đang mở & Quét mã QR bằng camera điện thoại hoặc mở mục Tự điểm danh, hệ thống lấy toạ độ, sinh viên chụp ảnh và gửi & Ghi điểm danh kèm khoảng cách và ảnh bằng chứng. Chưa đồng ý, thiếu hoặc sai mã, ngoài vùng, vị trí giả lập, khuôn mặt không khớp: mỗi trường hợp có thông báo riêng \\
UC16 Duyệt nghỉ có phép (lãnh đạo Khoa) & Sinh viên thuộc Khoa hoặc lớp do Khoa quản lý; bản ghi đang ở trạng thái Vắng & Mở phiên hoặc trang sinh viên, bấm Duyệt có phép, nhập lý do & Trạng thái thành Có phép, ghi vào nhật ký sửa, phụ huynh nhận thông báo cập nhật. Không có lý do thì không lưu; ngoài Khoa bị từ chối \\
UC24 Đồng ý hoặc rút đồng ý (sinh viên) & Đăng nhập bằng tài khoản sinh viên & Vào mục Khuôn mặt, đọc nội dung có số phiên bản, đánh dấu rồi xác nhận; hoặc bấm Rút đồng ý & Đồng ý được lưu kèm phiên bản và cách ghi nhận; rút đồng ý thì ảnh và vector khuôn mặt bị xoá ngay. Nội dung vừa đổi phiên bản thì phải đọc lại \\
Đặt lại mật khẩu (quản trị viên, Phòng Công tác sinh viên) & Tài khoản cần đặt lại thuộc phạm vi của người thao tác & Chọn Đặt lại mật khẩu ở dòng của người đó, nhập mật khẩu mới hai lần & Người đó bị đăng xuất khỏi mọi thiết bị, OTP còn hiệu lực bị huỷ, tài khoản đang khoá được mở. Hai lần nhập khác nhau hoặc vai trò không có quyền đều bị từ chối \\
\end{xltabular}}

Ở cấp Khoa, việc tách hai tác nhân xuất phát từ nguyên tắc \term{phân tách
nhiệm vụ} \cite{Sandhu1996RBAC}: người nhập dữ liệu không đồng thời là người phê duyệt dữ liệu đó.
Hình~\ref{fig:uckhoa} cho thấy hai tác nhân dùng chung nhóm ca sử dụng chỉ
xem, còn nhóm ca sử dụng làm thay đổi dữ liệu thì tách hẳn về hai phía.

\begin{figure}[H]
\centering
\begin{tikzpicture}[every node/.append style={align=center},
  uc/.style={ellipse, draw=daunavy, line width=0.7pt, fill=white, inner sep=0pt,
             minimum width=3.1cm, minimum height=0.9cm, font=\scriptsize},
  ucc/.style={uc, draw=daumaroon},
  lk/.style={line width=0.6pt, color=charcoal},
  actor/.style={line width=0.8pt}]
  % ranh gioi he thong
  \draw[line width=0.6pt, color=charcoal, rounded corners=3pt] (2.2,-4.3) rectangle (13.8,2.75);
  \node[anchor=north west, font=\scriptsize\bfseries] at (2.3,2.7) {FAMS, phạm vi một Khoa};
  % cot giua: dung chung, chi xem
  \node[ucc] (s1) at (8.0,1.9) {Xem sinh viên,\\ giảng viên, lớp};
  \node[ucc] (s2) at (8.0,0.8) {Xem điểm danh,\\ báo cáo, cảnh báo};
  % cot trai: van phong Khoa nhap lieu
  \node[uc] (v1) at (4.3,-0.5)  {Thêm, sửa\\ lớp học phần};
  \node[uc] (v2) at (4.3,-1.55) {Xếp lịch học,\\ chọn phòng};
  \node[uc] (v3) at (4.3,-2.6)  {Ghi danh\\ sinh viên};
  \node[uc] (v4) at (4.3,-3.65) {Phân công\\ giảng viên};
  % cot phai: lanh dao Khoa phe duyet, giam sat
  \node[uc] (l1) at (11.7,-0.5)  {Duyệt nghỉ\\ có phép};
  \node[uc] (l2) at (11.7,-1.55) {Xem nhật ký\\ của Khoa};
  % tac nhan
  \foreach \x in {1.0, 15.0}
    \draw[actor] (\x,-0.45) circle (0.16) (\x,-0.61)--(\x,-1.05) (\x-0.28,-0.75)--(\x+0.28,-0.75)
                 (\x,-1.05)--(\x-0.2,-1.4) (\x,-1.05)--(\x+0.2,-1.4);
  \node[font=\scriptsize] at (1.0,-1.85) {Văn phòng\\ Khoa};
  \node[font=\scriptsize] at (15.0,-1.85) {Lãnh đạo\\ Khoa};
  \foreach \u in {v1, v2, v3, v4, s1, s2} \draw[lk] (1.35,-0.75) -- (\u.west);
  \foreach \u in {l1, l2, s1, s2} \draw[lk] (14.65,-0.75) -- (\u.east);
\end{tikzpicture}
\caption{Ca sử dụng của hai tác nhân cấp Khoa: dùng chung nhóm chỉ xem, tách
riêng nhóm nhập liệu và nhóm phê duyệt.}
\label{fig:uckhoa}
\end{figure}

\subsection{Yêu cầu chức năng}

Nhóm chức năng cốt lõi gồm: quản lý danh mục và hồ sơ, trong đó có danh mục
phòng học; đồng ý sinh trắc học, đăng ký cùng duyệt dữ liệu khuôn mặt; vòng
đời phiên điểm danh gồm mở phiên, nhận diện, chiếu mã QR, xác nhận tay, tự
điểm danh và đóng phiên; sửa điểm danh có lưu vết lý do; thông báo điểm danh
cho phụ huynh; báo cáo chuyên cần, cảnh báo vắng và xuất Excel; thống kê tổng
hợp cho cấp quản lý; nhật ký hệ thống.

\subsection{Yêu cầu phi chức năng}

\begin{itemize}[leftmargin=1.6em, itemsep=2pt]
  \item \textbf{Hiệu năng}: nhận diện một khung hình phải đủ nhanh để giảng
        viên quét liên tục trong lớp, trên máy chủ chỉ có CPU.
  \item \textbf{An toàn}: mật khẩu băm bằng bcrypt \cite{Provos1999Bcrypt}; phiên đăng nhập thu hồi
        được; ảnh khuôn mặt và ảnh bằng chứng không truy cập được nếu không
        có quyền.
  \item \textbf{Riêng tư}: không có đồng ý sinh trắc học thì không có xử lý
        dữ liệu khuôn mặt; sinh viên xem được, rút đồng ý được và khi rút thì
        dữ liệu khuôn mặt bị xoá ngay.
  \item \textbf{Khả năng kiểm chứng}: mọi thay đổi trạng thái điểm danh đều
        lưu người thực hiện, thời điểm và lý do; mọi thao tác quản lý và mọi
        lần đăng nhập đều nằm trong nhật ký hệ thống.
  \item \textbf{Đa nền tảng}: cùng một mã nguồn chạy trên Web, Android, iOS.
\end{itemize}

\section{Kiến trúc tổng thể}

Hệ thống theo kiến trúc ba lớp \cite{Fowler2002PoEAA}, tách rõ ứng dụng người dùng, dịch vụ nghiệp
vụ và tầng lưu trữ (Hình~\ref{fig:kientruc}). Toàn bộ suy luận mô hình nằm ở
phía máy chủ, vì mô hình nặng khoảng 300\,MB và không thể phân phát xuống mọi
thiết bị của sinh viên.

\begin{figure}[H]
\centering
\begin{tikzpicture}[node distance=0.8cm and 0.8cm, every node/.append style={align=center}]
  \node (web)   [objbox] {Flutter Web\\ \scriptsize cán bộ, giảng viên};
  \node (mobi)  [objbox, below=of web] {Flutter Android, iOS\\ \scriptsize sinh viên, phụ huynh};
  \node (api)   [clsbox, right=of web, yshift=-1.05cm, minimum height=1.9cm]
        {FastAPI\\ \scriptsize \FsoEndpoint{} endpoint\\ \scriptsize xác thực, phân quyền};
  \node (svc)   [clsbox, right=of api, minimum height=1.9cm]
        {Lớp dịch vụ\\ \scriptsize nhận diện, chống giả mạo,\\ \scriptsize phiên, vị trí, thống kê};
  \node (db)    [membox, right=of svc, yshift=0.55cm] {PostgreSQL 16\\ \scriptsize pgvector, HNSW};
  \node (onnx)  [membox, right=of svc, yshift=-0.7cm] {ONNX Runtime\\ \scriptsize SCRFD, ArcFace,\\ \scriptsize MiniFASNetV2};
  \draw [flowarr] (web)  -- node[above, font=\scriptsize] {HTTPS} (api);
  \draw [flowarr] (mobi) -- (api);
  \draw [flowarr] (api)  -- (svc);
  \draw [flowarr] (svc)  -- (db);
  \draw [flowarr] (svc)  -- (onnx);
\end{tikzpicture}
\caption{Kiến trúc ba lớp của hệ thống FAMS.}
\label{fig:kientruc}
\end{figure}

Ba vòng lặp nền chạy song song với dịch vụ chính: một vòng dọn ảnh bằng
chứng quá hạn lưu trữ \FngayLuuAnh{} ngày, một vòng tự động đóng những phiên
điểm danh bị quên đóng, và một vòng gửi email thông báo cho phụ huynh. Cả ba
tắt được bằng tham số cấu hình để môi trường kiểm thử không chạm vào cơ sở dữ
liệu thật hay gửi thư ra ngoài.

\section{Thiết kế cơ sở dữ liệu}

Lược đồ gồm \FsoBang{} bảng, quản lý bằng \FsoMigration{} phiên bản di trú
Alembic. Bảng~\ref{tab:luocdo} liệt kê toàn bộ.

{\small
\begin{xltabular}{\textwidth}{@{}>{\raggedright\arraybackslash\leavevmode}p{5.0cm} X@{}}
\caption{Các bảng trong cơ sở dữ liệu.}
\label{tab:luocdo}\\
\toprule
Bảng & Nội dung \\
\midrule
\endfirsthead
\toprule
Bảng & Nội dung \\
\midrule
\endhead
\bottomrule
\endfoot
\BangDuLieu
\end{xltabular}}

{\small
\begin{xltabular}{\textwidth}{@{}>{\bfseries\raggedright\arraybackslash\leavevmode}p{2.4cm} >{\raggedright\arraybackslash\leavevmode}p{4.4cm} >{\raggedright\arraybackslash\leavevmode}X@{}}
\caption{Giải thích chi tiết Bảng~\ref{tab:luocdo}: các nhóm bảng và liên kết chính giữa chúng.}
\label{tab:luocdo-gt}\\
\toprule
Nhóm & Bảng & Liên kết chính \\
\midrule
\endfirsthead
\toprule
Nhóm & Bảng & Liên kết chính \\
\midrule
\endhead
\bottomrule
\endfoot
Tài khoản, đơn vị & \code{users}, \code{units}, \code{staff} & \code{staff} trỏ tới \code{users} và \code{units}: mỗi tài khoản cán bộ thuộc đúng một đơn vị \\
Hồ sơ & \code{students}, \code{lecturers}, \code{parents}, \code{parent\_students} & Ba bảng hồ sơ đều trỏ tới \code{users}; sinh viên và giảng viên trỏ tới \code{units} để biết thuộc Khoa nào; \code{parent\_students} nối phụ huynh với sinh viên \\
Đào tạo & \code{classes}, \code{enrollments}, \code{schedules}, \code{locations}, \code{rooms} & Lớp trỏ tới giảng viên và Khoa; \code{enrollments} nối sinh viên với lớp, mỗi cặp chỉ một lần; lịch học trỏ tới lớp, phòng và địa điểm; phòng thuộc một địa điểm \\
Điểm danh & \code{attendance\_sessions}, \code{attendance\_records}, \code{attendance\_audits} & Phiên trỏ tới lịch học, mỗi lịch một phiên một ngày; bản ghi trỏ tới phiên và sinh viên, mỗi sinh viên một bản ghi trong phiên; nhật ký sửa trỏ tới bản ghi và người sửa \\
Sinh trắc & \code{face\_embeddings} & Mỗi vector trỏ tới một sinh viên, kèm trạng thái duyệt \\
Bảo mật & \code{refresh\_tokens}, \code{password\_resets}, \code{revoked\_tokens} & Hai bảng đầu trỏ tới \code{users}; \code{revoked\_tokens} chỉ lưu mã định danh của access token bị thu hồi, không có khoá ngoại \\
Thông báo, nhật ký & \code{notifications}, \code{audit\_logs} & Thông báo bị xoá theo người nhận và để trống liên kết khi bản ghi điểm danh bị xoá; nhật ký giữ lại khi người thao tác hay Khoa bị xoá, chỉ để trống liên kết \\
\end{xltabular}}

Bốn quyết định thiết kế đáng chú ý:

\begin{enumerate}[leftmargin=1.8em, itemsep=3pt]
  \item \textbf{Vector đặc trưng nằm chung cơ sở dữ liệu quan hệ.} Cột
        \code{embedding} của bảng \code{face\_embeddings} có kiểu
        \code{Vector(\FchieuVector)} do pgvector cung cấp, kèm chỉ mục HNSW.
        Nhờ vậy một câu truy vấn duy nhất vừa lọc theo lớp học phần và trạng
        thái duyệt, vừa sắp xếp theo khoảng cách cosine.
  \item \textbf{Ràng buộc duy nhất đặt ở tầng cơ sở dữ liệu.} Mỗi lịch học
        chỉ có một phiên trong một ngày, mỗi sinh viên chỉ có một bản ghi
        trong một phiên, mỗi sinh viên chỉ ghi danh một lần vào một lớp.
        Những ràng buộc này được cơ sở dữ liệu bảo đảm \cite{Haerder1983Principles}, nên hai yêu cầu gửi
        đồng thời cũng không tạo được bản ghi trùng.
  \item \textbf{Mọi thao tác sửa điểm danh đều để lại vết.} Bảng
        \code{attendance\_audits} lưu trạng thái cũ, trạng thái mới, người
        sửa và lý do bắt buộc. Đây là điều kiện cần để kết quả chuyên cần có
        giá trị khi bị khiếu nại.
  \item \textbf{Nhật ký hệ thống tách khỏi nhật ký sửa điểm danh.} Bảng
        \code{audit\_logs} ghi mọi thao tác quản lý trên toàn hệ thống, kèm
        cột \code{faculty\_id} chốt lại Khoa liên quan ngay lúc ghi để lãnh
        đạo Khoa lọc được phần của mình. Thiết kế chi tiết ở
        Mục~\ref{sec:nhatky}.
\end{enumerate}

\section{Thiết kế giao diện lập trình}

Dịch vụ cung cấp \FsoEndpoint{} endpoint chia theo \FsoModuleApi{} nhóm chức
năng (Bảng~\ref{tab:api}). Toàn bộ đều yêu cầu access token, trừ đăng nhập và
luồng quên mật khẩu.

{\small
\begin{xltabular}{\textwidth}{@{}>{\raggedright\arraybackslash\leavevmode}p{3.1cm} c X@{}}
\caption{Các nhóm endpoint của dịch vụ.}
\label{tab:api}\\
\toprule
Tiền tố & Số lượng & Chức năng \\
\midrule
\endfirsthead
\toprule
Tiền tố & Số lượng & Chức năng \\
\midrule
\endhead
\bottomrule
\endfoot
\BangApi
\end{xltabular}}

{\small
\begin{xltabular}{\textwidth}{@{}>{\bfseries\raggedright\arraybackslash\leavevmode}p{3.4cm} >{\raggedright\arraybackslash\leavevmode}X@{}}
\caption{Giải thích chi tiết Bảng~\ref{tab:api}: quy ước chung của mọi endpoint.}
\label{tab:api-gt}\\
\toprule
Mục & Quy ước \\
\midrule
\endfirsthead
\toprule
Mục & Quy ước \\
\midrule
\endhead
\bottomrule
\endfoot
Tiền tố & Mỗi nhóm có tiền tố riêng như \code{/students}; khi chạy sau nginx, toàn bộ API nằm dưới \code{/api} \\
Số lượng & Số endpoint của nhóm, đếm tự động từ mã nguồn \\
Phương thức & GET để đọc; POST để tạo mới hoặc thực hiện một hành động như mở phiên, tự điểm danh, đặt lại mật khẩu; PATCH để sửa một phần; DELETE để xoá \\
Xác thực & Gửi kèm \code{Authorization: Bearer} cùng access token, trừ đăng nhập, làm mới token và quên mật khẩu bằng OTP; ảnh bằng chứng dùng đường dẫn được ký điện tử thay cho token \\
Phân trang & Danh sách lớn nhận \code{q}, \code{offset}, \code{limit} và trả tổng số bản ghi trong header \code{X-Total-Count} \\
Mã lỗi & 400 sai quy tắc nghiệp vụ; 401 chưa đăng nhập hoặc token hết hạn; 403 không đủ quyền hoặc ngoài phạm vi; 404 không tồn tại; 409 xung đột trạng thái, như nội dung đồng ý đã đổi phiên bản hay xoá lớp đang có sinh viên; 422 sai kiểu hoặc thiếu trường; 429 đăng nhập sai quá nhiều lần \\
\end{xltabular}}

Về xác thực, hệ thống dùng access token dạng JWT \cite{RFC7519} có thời hạn \FphutToken{}
phút, kèm refresh token xoay vòng \cite{RFC9700} thời hạn \FngayRefresh{} ngày cho ứng dụng
di động. Mỗi lần làm mới, refresh token cũ bị thu hồi; nếu một token đã thu
hồi được dùng lại, hệ thống coi là dấu hiệu bị đánh cắp và thu hồi toàn bộ
token của tài khoản đó.

Người quên mật khẩu có hai đường lấy lại. Đường thứ nhất là mã OTP \cite{NIST80063B4} gửi qua
email, người dùng tự làm. Đường thứ hai dành cho trường hợp email sai hoặc
không dùng được: cán bộ đặt lại mật khẩu cho người đó. Quản trị viên làm được
với mọi tài khoản, Phòng Công tác sinh viên làm được với sinh viên và phụ
huynh, tức đúng những hồ sơ phòng này quản lý. Mỗi lần đặt lại, hệ thống thu
hồi toàn bộ phiên đăng nhập, huỷ OTP còn hiệu lực và mở khoá nếu tài khoản
đang bị khoá vì nhập sai; nhật ký ghi rõ ai đã đổi cho ai nhưng không lưu giá
trị mật khẩu.

Ảnh khuôn mặt và ảnh bằng chứng không nằm trong thư mục tĩnh công khai. Dịch
vụ trả về \term{đường dẫn ảnh được ký điện tử}: máy chủ dùng khoá bí mật của
mình tính mã xác thực HMAC-SHA256 \cite{RFC2104} trên đường dẫn và mốc hết hạn
\FphutHetHanUrl{} phút, rồi gắn mã đó vào đường dẫn. Khi nhận yêu cầu xem ảnh,
máy chủ tính lại và so khớp; sửa một ký tự trong đường dẫn hay dùng sau mốc
hết hạn đều bị từ chối, nên một liên kết bị chia sẻ ra ngoài cũng tự mất hiệu
lực.

\section{Mô hình phân quyền}
\label{sec:phanquyen}

Phân quyền trong hệ thống có hai chiều. Chiều thứ nhất là \term{quyền theo
chức năng}, dựa trên mô hình phân quyền theo vai trò \cite{Sandhu1996RBAC}: vai trò được làm gì. Chiều thứ hai là \term{phạm vi dữ liệu}: vai
trò được thấy dữ liệu của ai. Chiều thứ hai mới là phần khó, vì cùng thao tác
xem danh sách sinh viên, Phòng Công tác sinh viên thấy toàn trường còn cán
bộ Khoa chỉ thấy sinh viên thuộc Khoa mình.

Đọc Bảng~\ref{tab:phanquyen} theo cột, mỗi phòng ban có một nhóm chức năng rõ ràng. Phòng Đào tạo nắm
phần tổ chức đào tạo toàn trường. Phòng Công tác sinh viên nắm hồ sơ người
học và mọi thứ liên quan tới dữ liệu sinh trắc. Văn phòng Khoa là cánh tay
nhập liệu của Khoa, làm cùng loại việc với Phòng Đào tạo nhưng chỉ trong Khoa
mình. Lãnh đạo Khoa không có quyền sửa dữ liệu gốc nào, chỉ có quyền xem và
quyền phê duyệt. Nhờ vậy người xếp lịch không tự duyệt đơn nghỉ, người duyệt
đơn không sửa được danh sách lớp, và thao tác của cả hai đều nằm trong nhật
ký mà lãnh đạo Khoa xem được.

{\footnotesize
\setlength{\tabcolsep}{4pt}
\begin{xltabular}{\textwidth}{@{}>{\raggedright\arraybackslash}X cccccc@{}}
\caption{Ma trận vai trò và chức năng của cán bộ. Ý nghĩa các ký hiệu ở
Bảng~\ref{tab:phanquyen-gt}.}
\label{tab:phanquyen}\\
\toprule
Chức năng & Quản trị & P.\,ĐT & P.\,CTSV & LĐ Khoa & VP Khoa & Ban ngành \\
\midrule
\endfirsthead
\toprule
Chức năng & Quản trị & P.\,ĐT & P.\,CTSV & LĐ Khoa & VP Khoa & Ban ngành \\
\midrule
\endhead
\bottomrule
\endfoot
Tài khoản cán bộ, đơn vị            & Sửa   & --  & --    & --     & --   & -- \\
Đặt lại mật khẩu giảng viên, cán bộ & Sửa   & --  & --    & --     & --   & -- \\
Đặt lại mật khẩu sinh viên, phụ huynh & Sửa & --  & Sửa   & --     & --   & -- \\
Lớp học phần, lịch, ghi danh,
  phân công giảng viên dạy lớp      & Sửa   & Sửa & Xem   & Xem*   & Sửa* & -- \\
Phòng học, địa điểm điểm danh       & Sửa   & Sửa & Xem   & Xem    & Xem  & -- \\
Hồ sơ giảng viên                    & Sửa   & Sửa & --    & Xem*   & Xem* & -- \\
Hồ sơ sinh viên, phụ huynh          & Sửa   & Xem & Sửa   & Xem*   & Xem* & -- \\
Ghi nhận đồng ý theo bản giấy       & Sửa   & --  & Sửa   & --     & --   & -- \\
Duyệt khuôn mặt                     & Sửa   & --  & Sửa   & --     & --   & -- \\
Điểm danh chi tiết và ảnh           & Sửa   & Xem & Xem   & Xem*   & Xem* & -- \\
Duyệt nghỉ có phép                  & Duyệt & --  & Duyệt & Duyệt* & --   & -- \\
Cảnh báo vắng, xuất Excel           & Xem   & Xem & Xem   & Xem*   & Xem* & -- \\
Thống kê tổng hợp                   & Xem   & Xem & Xem   & Xem*   & Xem* & Xem \\
Nhật ký hệ thống                    & Xem   & --  & --    & Xem*   & --   & -- \\
\end{xltabular}}

{\small
\begin{xltabular}{\textwidth}{@{}>{\bfseries\raggedright\arraybackslash\leavevmode}p{2.4cm} >{\raggedright\arraybackslash\leavevmode}X >{\raggedright\arraybackslash\leavevmode}p{5.0cm}@{}}
\caption{Giải thích chi tiết Bảng~\ref{tab:phanquyen}: ý nghĩa các ký hiệu.}
\label{tab:phanquyen-gt}\\
\toprule
Ký hiệu & Ý nghĩa & Ví dụ trong bảng \\
\midrule
\endfirsthead
\toprule
Ký hiệu & Ý nghĩa & Ví dụ trong bảng \\
\midrule
\endhead
\bottomrule
\endfoot
Sửa & Thêm, sửa, xoá được & Phòng Đào tạo với lớp học phần \\
Xem & Chỉ xem; không có nút sửa, gọi thẳng API sửa thì nhận lỗi 403 & Phòng Công tác sinh viên với lớp học phần \\
Duyệt & Chỉ được chuyển Vắng thành Có phép hoặc huỷ duyệt; không sửa được trạng thái khác & Phòng Công tác sinh viên, lãnh đạo Khoa với nghỉ có phép \\
Dấu sao & Quyền chỉ áp dụng trong Khoa của mình; dữ liệu Khoa khác không hiện ra & Văn phòng Khoa sửa lớp học phần của Khoa mình \\
Gạch ngang & Không có quyền; mục đó không có trong menu và API trả lỗi 403 & Ban ngành với hồ sơ sinh viên \\
Tên cột & P.\,ĐT là Phòng Đào tạo, P.\,CTSV là Phòng Công tác sinh viên, LĐ là lãnh đạo, VP là văn phòng & -- \\
\end{xltabular}}


Ba vai trò còn lại có phạm vi hẹp và cố định: giảng viên thao tác trên lớp
mình dạy; sinh viên chỉ truy cập dữ liệu của chính mình; phụ huynh chỉ xem
chuyên cần của con đã được cán bộ liên kết.

\begin{ghinho}
Phạm vi dữ liệu được cài đặt tập trung tại một mô-đun phân quyền duy nhất
dưới dạng bộ lọc truy vấn, thay vì rải điều kiện \code{if} ở từng endpoint.
Nhờ vậy khi thêm một vai trò mới, chỉ cần sửa một nơi, và bộ kiểm thử phân
quyền kiểm soát được toàn bộ ma trận trên.
\end{ghinho}

Vai trò của cán bộ không nhập tay mà suy ra từ loại đơn vị khi tạo tài khoản:
loại Phòng Đào tạo cho ra vai trò Phòng Đào tạo, và tương tự với các loại còn
lại. Riêng đơn vị loại Khoa có hai vai trò, quản trị viên chọn văn phòng Khoa
hoặc lãnh đạo Khoa; chọn lãnh đạo Khoa cho một đơn vị không phải Khoa sẽ bị
từ chối. Cách này loại bỏ hẳn tình trạng gán nhầm vai trò không khớp với đơn
vị.

\section{Thiết kế luồng điểm danh}

\subsection{Luồng đăng ký khuôn mặt}

Trước khi được điểm danh tự động, sinh viên phải đồng ý sử dụng dữ liệu sinh
trắc học (Mục~\ref{sec:dongy}) rồi mới đăng ký được khuôn mặt. Ảnh gửi lên đi
qua bốn cửa: phải có đúng một khuôn mặt trong ảnh; ảnh phải đạt ngưỡng
về kích thước khuôn mặt, độ sáng và độ nét; phải vượt kiểm tra chống giả mạo;
và vector thu được không được quá giống vector đã duyệt của một sinh viên
khác, cụ thể là độ tương tự không vượt \FnguongTrung. Cửa cuối cùng nhằm chặn
tình huống một người đăng ký khuôn mặt của mình dưới hai danh tính.

Vector qua được cả bốn cửa vẫn ở trạng thái \term{chờ duyệt} và chưa dùng để
nhận diện. Cán bộ Phòng Công tác sinh viên đối chiếu ảnh rồi duyệt hoặc từ
chối; ngay sau thao tác đó, ảnh khuôn mặt tạm bị xoá, chỉ vector ở lại.

\subsection{Các trạng thái điểm danh}

Một bản ghi điểm danh mang một trong năm trạng thái: \term{có mặt} khi được
ghi nhận trong khung giờ đầu tiết; \term{muộn} khi được ghi nhận sau giờ bắt
đầu quá \FphutMuon{} phút; \term{vắng} do hệ thống sinh ra lúc đóng phiên;
\term{xác nhận tay} khi giảng viên tự đánh dấu; và \term{có phép} khi Phòng
Công tác sinh viên hoặc lãnh đạo Khoa duyệt đơn nghỉ. Chỉ trạng thái vắng
không phép mới tính vào tỉ lệ cảnh báo cấm thi.

\subsection{Vòng đời một phiên điểm danh}

Phiên điểm danh là trung tâm của nghiệp vụ. Hình~\ref{fig:vongdoiphien} mô tả
vòng đời của nó.

\begin{figure}[H]
\centering
\begin{tikzpicture}[node distance=1.0cm and 1.0cm, every node/.append style={align=center}]
  \node (mo)    [startstop] {Mở phiên\\ \scriptsize sớm nhất \FphutMoSom{} phút trước tiết};
  \node (nhan)  [process, right=of mo, xshift=0.5cm] {Nhận điểm danh\\ \scriptsize camera, tự chụp, tay};
  \node (han)   [decision, right=of nhan, xshift=0.8cm] {Quá hạn nhận?};
  \node (dong)  [process, below=of han] {Đóng phiên};
  \node (vang)  [startstop, left=of dong, xshift=-0.8cm] {Sinh bản ghi vắng\\ \scriptsize và cảnh báo};
  \draw [flowarr] (mo)   -- (nhan);
  \draw [flowarr] (nhan) -- (han);
  \draw [flowarr] (han)  -- node[right, font=\scriptsize] {có} (dong);
  \draw [flowarr] (han.north) to[out=60, in=20, looseness=2.2]
        node[right, font=\scriptsize] {chưa} (nhan.north east);
  \draw [flowarr] (dong) -- (vang);
\end{tikzpicture}
\caption{Vòng đời của một phiên điểm danh.}
\label{fig:vongdoiphien}
\end{figure}

Phiên chỉ mở được trong khoảng từ \FphutMoSom{} phút trước giờ bắt đầu tiết
tới giờ kết thúc. Sau giờ kết thúc \FphutAnHan{} phút, hệ thống ngừng nhận
điểm danh và tự động đóng phiên nếu giảng viên quên. Khi đóng, mọi sinh viên
đã ghi danh mà chưa có bản ghi sẽ được ghi nhận vắng, đồng thời hệ thống tính
lại tỉ lệ vắng và phát cảnh báo cho những sinh viên vừa vượt ngưỡng
\FtiLeCanhBao\%.

\subsection{Nhận diện một khung hình}

Giải thuật~\ref{alg:nhandien} mô tả xử lý phía máy chủ khi giảng viên gửi lên
một khung hình từ camera.

\begin{algorithm}[H]
\caption{Nhận diện điểm danh từ một khung hình của lớp}
\label{alg:nhandien}
\hspace*{\algorithmicindent}\textbf{Đầu vào:} ảnh $I$, mã lớp $c$, ngưỡng $\tau = \FnguongNhanDien$\\
\hspace*{\algorithmicindent}\textbf{Đầu ra:} danh sách danh tính nhận được kèm độ tin cậy
\begin{algorithmic}[1]
\Function{NhanDienKhungHinh}{$I$, $c$}
  \State $F \gets$ \Call{PhatHienKhuonMat}{$I$}
  \State $R \gets \emptyset$
  \ForAll{$f \in F$}
    \State $(that, diem) \gets$ \Call{KiemTraSuSong}{$I$, $f$}
    \State $e \gets$ \Call{TrichDacTrung}{$I$, $f$}
    \State $(s, d) \gets$ láng giềng gần nhất của $e$ trong các vector
           \textbf{đã duyệt} của sinh viên thuộc lớp $c$ \textbf{và đã đồng
           ý} sinh trắc học \label{ln:lanhgieng}
    \State $conf \gets 1 - d$
    \If{$conf \geq \tau$ \textbf{và} $that$}
      \State $R \gets R \cup \{(s, conf)\}$
    \Else
      \State $R \gets R \cup \{(\textsc{Null}, conf)\}$
           \Comment{ghi nhận để giảng viên biết có người chưa nhận ra}
    \EndIf
  \EndFor
  \State \Return $R$
\EndFunction
\end{algorithmic}
\end{algorithm}

Ba chi tiết quan trọng nằm ở dòng~\ref{ln:lanhgieng}. Thứ nhất, phép tìm
láng giềng chỉ chạy trên sinh viên thuộc lớp đang điểm danh, nên một sinh viên
lớp khác đi ngang qua camera không bao giờ bị điểm danh nhầm. Thứ hai, chỉ
những vector đã được cán bộ duyệt mới tham gia so khớp; vector đang chờ duyệt
không đủ tin cậy để làm căn cứ. Thứ ba, sinh viên chưa đồng ý hoặc đã rút đồng
ý không bao giờ được nhận ra, kể cả khi vì lý do nào đó vector cũ còn sót lại
trong cơ sở dữ liệu.

\subsection{Sinh viên tự điểm danh}

Ngoài luồng giảng viên quét camera, sinh viên có thể tự điểm danh bằng ảnh tự
chụp kèm toạ độ GPS, và nếu giảng viên yêu cầu thì kèm cả mã QR đang chiếu
trong phòng. Luồng này rủi ro hơn hẳn, nên bị siết bằng các điều kiện mô tả
trong Giải thuật~\ref{alg:tudiemdanh}.

\begin{algorithm}[H]
\caption{Sinh viên tự điểm danh bằng ảnh tự chụp và vị trí}
\label{alg:tudiemdanh}
\hspace*{\algorithmicindent}\textbf{Đầu vào:} ảnh $I$, sinh viên $s$, phiên $p$, mã QR $q$, toạ độ $(lat, lon)$, sai số $a$\\
\hspace*{\algorithmicindent}\textbf{Đầu ra:} chấp nhận hoặc lý do từ chối
\begin{algorithmic}[1]
\Function{TuDiemDanh}{$I$, $s$, $p$, $q$, $lat$, $lon$, $a$}
  \If{$s$ chưa ghi danh vào lớp của $p$ \textbf{hoặc} $p$ đã quá hạn nhận}
    \State \Return \textsc{TuChoi}
  \EndIf
  \If{$s$ chưa đồng ý sử dụng dữ liệu sinh trắc học} \label{ln:dongy}
    \State \Return \textsc{CanDongY} \Comment{giảng viên vẫn xác nhận tay được}
  \EndIf
  \If{$p$ yêu cầu mã QR \textbf{và} $q$ không khớp khung giờ hiện tại hay liền trước} \label{ln:qr}
    \State \Return \textsc{TuChoi}
  \EndIf
  \State $L \gets$ vùng của phòng học, nếu không có thì khu vực chứa phòng, rồi tới địa điểm mặc định
  \If{$L \neq \textsc{Null}$}
    \If{vị trí bị giả lập \textbf{hoặc} $a > \FsaiSoGps$ mét}
      \State \Return \textsc{TuChoi}
    \EndIf
    \State $dist \gets$ \Call{Haversine}{$lat$, $lon$, $L.lat$, $L.lon$}
    \If{$dist - a > L.banKinh$}
      \State \Return \textsc{TuChoi} \Comment{đang ở ngoài phạm vi cho phép}
    \EndIf
  \EndIf
  \State $(trangThai, conf) \gets$ so khớp $I$ \textbf{chỉ với} vector đã
         duyệt của chính $s$ \label{ln:chinhminh}
  \If{$trangThai \neq \textsc{Khop}$}
    \State \Return \textsc{TuChoi}
  \EndIf
  \State ghi bản ghi điểm danh với nguồn \textsc{selfie}, kèm $dist$ và ảnh bằng chứng
  \State \Return \textsc{ChapNhan}
\EndFunction
\end{algorithmic}
\end{algorithm}

Điểm khác biệt cốt yếu so với luồng camera nằm ở dòng~\ref{ln:chinhminh}:
ảnh tự chụp chỉ được so khớp với vector của chính sinh viên đang đăng nhập,
chứ không so với cả lớp. Nếu so với cả lớp, một sinh viên có thể chụp ảnh bạn
mình và hệ thống vẫn tìm ra một danh tính hợp lệ. Dòng~\ref{ln:dongy} trả về
một kết quả riêng thay vì từ chối chung chung, để ứng dụng đưa sinh viên tới
màn hình đọc nội dung đồng ý thay vì để sinh viên chụp lại mãi.

Khoảng cách tới địa điểm được tính bằng công thức Haversine trên mặt cầu \cite{Sinnott1984Haversine}:

\begin{equation}
\label{eq:haversine}
d = 2R \arcsin \sqrt{ \sin^2\!\left(\frac{\Delta\varphi}{2}\right)
+ \cos\varphi_1 \cos\varphi_2 \sin^2\!\left(\frac{\Delta\lambda}{2}\right) }
\end{equation}

trong đó $\varphi$ là vĩ độ, $\lambda$ là kinh độ và $R$ là bán kính Trái
Đất. Vì GPS trong nhà có sai số hàng chục mét \cite{Zandbergen2011Positional}, hệ thống trừ đi sai số do thiết
bị báo lên trước khi so với bán kính.

Vùng đối chiếu $L$ được chọn từ hẹp tới rộng. Nếu phòng học của lịch có toạ
độ riêng thì dùng toạ độ đó, với bán kính không nhỏ hơn \FbanKinhPhong\,m để
không chặn oan sinh viên đang ngồi đúng lớp. Nếu phòng chưa có toạ độ thì dùng
khu vực chứa phòng, sau cùng mới tới địa điểm mặc định của trường. Danh mục
phòng học còn giúp chặn xếp hai lớp vào cùng một phòng trong cùng khung giờ
ngay khi nhập lịch.

\subsection{Mã QR điểm danh xoay vòng}

Toạ độ GPS do thiết bị gửi lên nên vẫn có thể bị làm giả. Với những buổi cần
chắc chắn, giảng viên bật yêu cầu mã QR khi mở phiên và chiếu mã lên màn hình
trong phòng \cite{ISO18004}. Mã được ký điện tử bằng HMAC \cite{RFC2104} trên mã phiên và số thứ tự khung thời gian
\FgiayDoiQr{} giây, nên tự đổi sau mỗi \FgiayDoiQr{} giây và chỉ máy chủ sinh
ra được, theo cùng ý tưởng với mật khẩu một lần dựa trên thời gian
\cite{RFC6238}. Máy chủ chấp nhận mã của khung hiện tại và khung liền trước, tức một
mã sống không quá hai khung. Ảnh chụp mã gửi cho bạn ở nhà vì thế hết hạn trước
khi kịp dùng. Kèm theo mã QR là một mã sáu chữ số sinh cùng cách, để sinh viên
có máy ảnh hỏng vẫn gõ tay được.

Ảnh QR chứa một đường link tới trang tự điểm danh của đúng buổi, kèm sẵn mã,
nên sinh viên quét bằng camera thường của điện thoại là vào thẳng trang đó mà
không phải mở ứng dụng rồi tìm mục quét mã. Nếu chưa đăng nhập, sinh viên đăng
nhập xong sẽ quay lại đúng trang này. Trong lúc mã đang được chiếu, màn hình
quét camera của giảng viên tự dừng nhận diện để webcam không ghi điểm danh cho
người ngồi trước máy.

\subsection{Thông báo cho phụ huynh}

Mỗi khi bản ghi điểm danh của sinh viên được tạo hoặc đổi trạng thái, kể cả
khi đóng phiên sinh ra bản ghi vắng, hệ thống tạo thông báo cho từng phụ huynh
đã liên kết ngay trong cùng giao dịch ghi điểm danh, nên không có bản ghi nào
thiếu thông báo. Thông báo nằm trong hộp thư của ứng dụng; nếu phụ huynh bật
nhận email thì thông báo còn được đưa vào hàng đợi, một vòng lặp nền gửi đi, mỗi
thư được thử tối đa \FsoLanGuiEmail{} lần nếu máy chủ thư lỗi. Phụ huynh tự
chọn có nhận thông báo có mặt hay không; thông báo vắng, đi muộn và các lần
sửa kết quả thì luôn được gửi.

\section{Các lớp phòng vệ chống gian lận}
\label{sec:baolop}

Không lớp phòng vệ nào là tuyệt đối, nên hệ thống xếp chồng sáu lớp, theo tinh
thần nguyên tắc tách đặc quyền \en{(separation of privilege)}
\cite{Saltzer1975Protection}. Một
hành vi gian lận phải vượt qua toàn bộ mới thành công.

\begin{table}[H]
\centering
\caption{Sáu lớp phòng vệ và kiểu tấn công mà mỗi lớp ngăn chặn.}
\label{tab:phongve}
\small
\begin{tabularx}{\textwidth}{@{}>{\raggedright\arraybackslash}p{4.2cm} X@{}}
\toprule
Lớp phòng vệ & Ngăn kiểu gian lận \\
\midrule
Kiểm tra chất lượng ảnh & Ảnh quá mờ, quá tối hoặc khuôn mặt quá nhỏ, vốn là
  điều kiện dễ khiến mô hình nhận nhầm \\
Chống giả mạo MiniFASNetV2 & Giơ ảnh in hoặc ảnh trên màn hình điện thoại
  trước camera \\
Xác nhận \FsoKhungXacNhan{} khung hình & Một khung hình may mắn vượt ngưỡng
  do góc chụp hoặc ánh sáng bất lợi \\
So khớp theo phạm vi & Nhận nhầm người không thuộc lớp; ở luồng tự điểm danh
  là chụp hộ bạn cùng lớp \\
Đối chiếu GPS theo phòng học & Điểm danh từ nhà hoặc từ toà nhà khác \\
Mã QR xoay vòng \FgiayDoiQr{} giây và quyền của giảng viên & Chụp mã gửi cho
  bạn ở xa; mọi trường hợp còn lại đều có thể bị giảng viên sửa lại kèm lý do \\
\bottomrule
\end{tabularx}
\end{table}

{\small
\begin{xltabular}{\textwidth}{@{}>{\bfseries\raggedright\arraybackslash\leavevmode}p{2.6cm} >{\raggedright\arraybackslash\leavevmode}p{5.2cm} >{\raggedright\arraybackslash\leavevmode}X@{}}
\caption{Giải thích chi tiết Bảng~\ref{tab:phongve}: mỗi lớp được cài đặt ở đâu và xử lý thế nào khi không vượt qua.}
\label{tab:phongve-gt}\\
\toprule
Lớp phòng vệ & Cài đặt và tham số & Khi không vượt qua \\
\midrule
\endfirsthead
\toprule
Lớp phòng vệ & Cài đặt và tham số & Khi không vượt qua \\
\midrule
\endhead
\bottomrule
\endfoot
Chất lượng ảnh & \code{face\_service.py}, áp dụng khi đăng ký khuôn mặt: cạnh mặt tối thiểu 112 điểm ảnh, độ sáng trung bình tối thiểu 40, độ nét theo phương sai Laplace tối thiểu 60 & Ảnh đăng ký bị từ chối kèm lý do, sinh viên chụp lại \\
Chống giả mạo & \code{antispoof\_service.py}, mô hình MiniFASNetV2; điểm mặt thật tối thiểu \FnguongThat & Khung hình bị đánh dấu nghi giả mạo, không ghi điểm danh \\
Xác nhận đa khung hình & Endpoint nhận diện của phiên trong \code{sessions.py}; cần \FsoKhungXacNhan{} lần nhận ra & Chưa tạo bản ghi cho tới khi đủ số lần \\
So khớp theo phạm vi & \code{face\_service.py}: chỉ xét sinh viên ghi danh của lớp, khuôn mặt đã duyệt, đã đồng ý; ngưỡng \FnguongNhanDien & Không trả về danh tính nào \\
GPS theo phòng học & \code{geo\_service.py}: sai số tối đa \FsaiSoGps{}\,m, bán kính phòng tối thiểu \FbanKinhPhong{}\,m & Từ chối với lý do riêng: ngoài vùng, sai số lớn hoặc vị trí giả lập \\
Mã QR và quyền giảng viên & \code{checkin\_code\_service.py}: đổi mã mỗi \FgiayDoiQr{} giây, nhận mã của khung hiện tại và khung liền trước & Từ chối khi thiếu hoặc sai mã; mọi trường hợp lọt qua vẫn có thể bị giảng viên sửa kèm lý do \\
\end{xltabular}}

\section{Thiết kế bảo vệ dữ liệu sinh trắc}

\subsection{Đồng ý sử dụng dữ liệu sinh trắc học là bắt buộc}
\label{sec:dongy}

Dữ liệu khuôn mặt là dữ liệu cá nhân nhạy cảm \cite{Luat912025, ND3562025}, nên hệ thống coi sự đồng ý của
sinh viên là điều kiện bắt buộc, được kiểm tra ở mọi điểm xử lý chứ không chỉ
hiện một hộp thoại cho có:

\begin{itemize}[leftmargin=1.6em, itemsep=3pt]
  \item \textbf{Đồng ý với một nội dung cụ thể.} Nội dung đồng ý nằm ở máy chủ
        và có phiên bản, hiện tại là \code{\FphienBanDongY}. Nội dung nêu rõ thu
        thập gì, để làm gì, giữ bao lâu, ai được xem và cách rút đồng ý. Khi
        xác nhận, ứng dụng gửi kèm phiên bản sinh viên đã đọc; nếu nội dung
        vừa đổi, máy chủ từ chối và ứng dụng hiện lại nội dung mới. Khi nhà
        trường sửa nội dung, sinh viên đã đồng ý bản cũ được nhắc đọc và xác
        nhận lại.
  \item \textbf{Hai cách ghi nhận, không có cách đồng ý thay.} Sinh viên tự
        tick và xác nhận trong ứng dụng; hoặc sinh viên ký bản giấy, cán bộ
        Phòng Công tác sinh viên nhập lại và hệ thống lưu tên người nhập. Cán
        bộ không thể đồng ý thay khi không có bản giấy.
  \item \textbf{Chưa đồng ý thì không xử lý.} Không đăng ký được khuôn mặt,
        không được nhận ra khi giảng viên quét camera, và không tự điểm danh
        được. Quyền học tập không bị ảnh hưởng: giảng viên vẫn xác nhận có
        mặt trực tiếp.
  \item \textbf{Rút đồng ý là xoá ngay.} Sinh viên tự rút trong ứng dụng, hoặc
        Phòng Công tác sinh viên làm thay khi sinh viên nộp đơn; toàn bộ ảnh và
        vector khuôn mặt bị xoá trong cùng thao tác. Mỗi lần đồng ý, rút đồng
        ý đều được ghi vào nhật ký hệ thống.
\end{itemize}

\subsection{Tối thiểu hoá và quyền kiểm soát dữ liệu}

Hai nguyên tắc còn lại chi phối cách hệ thống lưu dữ liệu khuôn mặt, tương ứng
với nguyên tắc tối thiểu hoá dữ liệu và quyền của chủ thể dữ liệu trong
\cite{GDPR2016}:

\begin{itemize}[leftmargin=1.6em, itemsep=3pt]
  \item \textbf{Chỉ giữ thứ tối thiểu.} Thứ lưu lâu dài là vector đặc trưng,
        không phải ảnh. Ảnh khuôn mặt chỉ được giữ tạm để cán bộ đối chiếu
        khi duyệt và bị xoá ngay sau khi duyệt hoặc từ chối. Ảnh bằng chứng
        điểm danh tự xoá sau \FngayLuuAnh{} ngày.
  \item \textbf{Người dùng kiểm soát được dữ liệu của mình.} Sinh viên xem
        được danh sách dữ liệu khuôn mặt đã đăng ký của mình và rút đồng ý,
        tức xoá toàn bộ, bất cứ lúc nào.
\end{itemize}

\begin{luuy}
Vector đặc trưng không phải là dạng ẩn danh tuyệt đối: đã có công trình dựng
lại được ảnh khuôn mặt từ vector đặc trưng \cite{Mai2019Reconstruction}. Vì vậy hệ thống vẫn
coi bảng \code{face\_embeddings} là dữ liệu nhạy cảm: không endpoint nào trả
vector ra ngoài, và mọi truy cập đều đi qua lớp phân quyền.
\end{luuy}

\section{Nhật ký hệ thống}
\label{sec:nhatky}

Nhật ký sửa điểm danh ở bảng \code{attendance\_audits} chỉ trả lời được ai
đã sửa kết quả điểm danh. Để trả lời câu hỏi rộng hơn là ai đã làm gì trên hệ
thống, bảng \code{audit\_logs} ghi lại ba nhóm sự kiện, theo hướng dẫn quản lý nhật ký an toàn thông tin \cite{Kent2006NIST}:

\begin{itemize}[leftmargin=1.6em, itemsep=3pt]
  \item \textbf{Đăng nhập và đăng xuất}, gồm cả lần đăng nhập sai và lần thử
        khi tài khoản đang bị khoá tạm, kèm địa chỉ IP.
  \item \textbf{Mọi thao tác thêm, sửa, xoá} trên tài khoản, đơn vị, cán bộ,
        sinh viên, giảng viên, phụ huynh, lớp học phần, ghi danh, lịch học,
        phòng học, địa điểm, dữ liệu khuôn mặt, phiên và kết quả điểm danh.
        Một số thay đổi mang tên riêng cho dễ đọc như đổi mật khẩu, đồng ý
        hay rút đồng ý sinh trắc học, duyệt khuôn mặt, mở và đóng phiên.
  \item \textbf{Xuất báo cáo Excel}, vì đây là lúc dữ liệu cá nhân rời khỏi
        hệ thống.
\end{itemize}

Mỗi dòng nhật ký lưu người thao tác, vai trò, thời điểm, hành động, đối
tượng, kết quả và phần thay đổi dưới dạng cặp giá trị cũ, giá trị mới của
từng trường.

\begin{ghinho}
Nhật ký được ghi tự động ở tầng ánh xạ đối tượng: một hàm lắng nghe sự kiện
\code{after\_flush} của SQLAlchemy đọc mọi đối tượng vừa thêm, sửa, xoá trong
giao dịch và ghi nhật ký vào chính giao dịch đó. Vì vậy không endpoint nào
quên ghi được, và nếu giao dịch bị huỷ thì dòng nhật ký cũng huỷ theo, không
có nhật ký cho thao tác chưa từng xảy ra.
\end{ghinho}

Ba chi tiết giữ cho nhật ký dùng được trong thực tế. Thứ nhất, các trường
nhạy cảm như mật khẩu đã băm, vector khuôn mặt hay mã OTP chỉ được ghi là
\emph{đã thay đổi}, không bao giờ ghi giá trị. Thứ hai, việc tạo bản ghi điểm
danh khi sinh viên được nhận diện là luồng bình thường, xảy ra hàng nghìn lần
mỗi ngày, nên không ghi; chỉ những lần có người sửa kết quả mới được ghi.
Thứ ba, Khoa liên quan được chốt ngay lúc ghi, nên lãnh đạo Khoa chỉ thấy
thao tác trên dữ liệu thuộc Khoa mình, và kết quả không đổi kể cả khi đối
tượng sau này bị xoá hoặc chuyển Khoa.

% ---------------------------------------------------------------------
%  CHƯƠNG 4 — CÀI ĐẶT HỆ THỐNG
% ---------------------------------------------------------------------
\chapter{Cài đặt hệ thống}
\label{ch:caidat}

\begin{muctieu}
Chương này mô tả môi trường phát triển, cách tổ chức mã nguồn, cài đặt của
ba thành phần cốt lõi gồm nhận diện khuôn mặt, đối chiếu vị trí và phân
quyền theo phạm vi, cách cài đặt đồng ý sinh trắc học và nhật ký hệ thống,
cùng cách ứng dụng đa nền tảng và quy trình triển khai.
\end{muctieu}

\section{Môi trường phát triển}

\begin{table}[H]
\centering
\caption{Môi trường phát triển và triển khai.}
\label{tab:moitruong}
\small
\begin{tabularx}{\textwidth}{@{}>{\bfseries}p{4.2cm} X@{}}
\toprule
Thành phần & Phiên bản và ghi chú \\
\midrule
Ngôn ngữ phía máy chủ & Python 3.10 \\
Khung dịch vụ & FastAPI, uvicorn, SQLAlchemy 2.0, Alembic \\
Cơ sở dữ liệu & PostgreSQL 16 kèm phần mở rộng pgvector \\
Thư viện thị giác & InsightFace \code{buffalo\_l}, OpenCV, ONNX Runtime bản CPU \\
Ứng dụng & Flutter (Dio, go\_router, provider, camera, geolocator, fl\_chart) \\
Đóng gói & Docker Compose ba dịch vụ: cơ sở dữ liệu, máy chủ, giao diện \\
Kiểm thử & pytest chạy trên PostgreSQL thật; \code{flutter test} cho ứng dụng \\
Máy thử nghiệm & Intel Core i5-1135G7 (4 nhân, 8 luồng), RAM 16\,GB, không dùng GPU \\
\bottomrule
\end{tabularx}
\end{table}

{\small
\begin{xltabular}{\textwidth}{@{}>{\bfseries\raggedright\arraybackslash\leavevmode}p{3.4cm} >{\raggedright\arraybackslash\leavevmode}X@{}}
\caption{Giải thích chi tiết Bảng~\ref{tab:moitruong}: vai trò của từng thành phần trong FAMS.}
\label{tab:moitruong-gt}\\
\toprule
Thành phần & Dùng để làm gì \\
\midrule
\endfirsthead
\toprule
Thành phần & Dùng để làm gì \\
\midrule
\endhead
\bottomrule
\endfoot
uvicorn & Máy chủ chạy ứng dụng FastAPI \\
SQLAlchemy 2.0 & Ánh xạ bảng thành lớp Python; sự kiện của nó được dùng để ghi nhật ký hệ thống và sinh thông báo phụ huynh \\
Alembic & Di trú lược đồ cơ sở dữ liệu, hiện có \FsoMigration{} phiên bản \\
pgvector & Kiểu \code{Vector(\FchieuVector)}, toán tử khoảng cách cosine và chỉ mục HNSW \\
InsightFace \code{buffalo\_l} & Bộ năm mô hình: SCRFD phát hiện khuôn mặt, ArcFace trích đặc trưng, và ba mô hình phụ (điểm mốc 2D, điểm mốc 3D, tuổi và giới tính) mà hệ thống không dùng tới, xem Mục~\ref{sec:hieunang} \\
OpenCV & Giải mã ảnh, cắt vùng khuôn mặt, đo độ sáng và độ nét \\
ONNX Runtime & Chạy các mô hình trên CPU \\
Gói Flutter & Dio gọi API và tự làm mới token; go\_router định tuyến và chặn trang theo vai trò; provider giữ trạng thái đăng nhập; camera, geolocator, fl\_chart cho camera, định vị và biểu đồ \\
Máy thử nghiệm & Windows 11, Python 3.10, ONNX Runtime 1.20 bản CPU; mọi số đo ở Chương~\ref{ch:ketqua} lấy trên máy này \\
\end{xltabular}}

Không cần GPU. Toàn bộ mô hình chạy qua ONNX Runtime trên CPU, đây là ràng
buộc tự đặt ra từ đầu, vì một hệ thống đòi hỏi máy chủ có GPU sẽ khó được
nhà trường triển khai thật.

\section{Tổ chức mã nguồn}

Mã nguồn chia hai phần lớn, tổng cộng khoảng \FsoDongPython{} dòng Python
cho phía máy chủ và \FsoDongDart{} dòng Dart cho ứng dụng.

\begin{lstlisting}[caption={Cấu trúc thư mục của kho mã nguồn.}, label={lst:caytm}]
FAMS/
|-- backend/
|   |-- alembic/versions/   # database migrations
|   |-- seed.py             # sample data
|   |-- tests/              # pytest suite
|   `-- app/
|       |-- api/            # HTTP endpoints, grouped by resource
|       |-- core/           # config, security, authz, signed media URLs
|       |-- db/             # ORM models and session factory
|       |-- services/       # face, antispoof, session, geo, audit, notify
|       `-- schemas.py      # request and response models
|-- mobile/lib/             # Flutter app for Web, Android, iOS
|   |-- core/               # api client, auth controller, role matrix
|   |-- features/           # one folder per role group
|   `-- widgets/            # shared UI components
`-- baocao/                 # this report and its generator
\end{lstlisting}

Tầng \code{api} chỉ nhận yêu cầu, kiểm tra quyền rồi gọi xuống tầng
\code{services}. Mọi phần liên quan tới mô hình nằm gọn trong
\code{services}, nên có thể đổi mô hình nhận diện mà không phải chạm vào
tầng endpoint.

\section{Cài đặt nhận diện khuôn mặt}

Mô hình được nạp theo kiểu trì hoãn: lần gọi đầu tiên mới tải bộ
\code{buffalo\_l}, những lần sau dùng lại phiên suy luận đã có trong bộ nhớ.
Nhờ đó máy chủ khởi động nhanh và tiến trình kiểm thử không phải tải mô hình
khi không cần.

Mã nguồn~\ref{lst:nhandien} là hàm xử lý một khung hình do giảng viên gửi
lên. Đây là hiện thực của Giải thuật~\ref{alg:nhandien}.

\lstinputlisting[language=Python, caption={Nhận diện toàn bộ khuôn mặt trong
một khung hình, giới hạn theo lớp học phần (đã lược chú thích).},
label={lst:nhandien}]{code/nhan-dien-khung-hinh.py}

Ba điểm đáng chú ý trong đoạn mã trên:

\begin{itemize}[leftmargin=1.6em, itemsep=3pt]
  \item \textbf{Phép so khớp do cơ sở dữ liệu thực hiện.} Biểu thức
        \code{cosine\_distance} sinh ra toán tử của pgvector, nên việc sắp
        xếp theo khoảng cách chạy ngay trong PostgreSQL và tận dụng chỉ mục
        HNSW, thay vì tải toàn bộ vector về máy chủ ứng dụng rồi duyệt.
  \item \textbf{Điều kiện lọc nằm cùng câu truy vấn.} Phép nối với bảng ghi
        danh, điều kiện trạng thái \code{approved} và điều kiện
        \code{consent\_given} khiến không gian tìm kiếm chỉ còn sinh viên của
        đúng lớp đó, chỉ những khuôn mặt đã được duyệt và chỉ những sinh viên
        đang đồng ý sử dụng dữ liệu sinh trắc học.
  \item \textbf{Khuôn mặt nghi giả mạo vẫn được nhận diện.} Hàm không bỏ qua
        khung hình nghi giả mạo mà vẫn trả về danh tính kèm cờ báo. Quyết
        định không ghi nhận điểm danh thuộc về tầng endpoint, nhờ vậy giảng
        viên biết ai đang cố gian lận thay vì chỉ thấy hệ thống im lặng.
\end{itemize}

Việc ghi nhận điểm danh còn phải qua bộ đếm xác nhận đa khung hình: một sinh
viên phải được nhận diện đủ \FsoKhungXacNhan{} lần trong cùng một phiên thì
bản ghi mới được tạo. Bộ đếm lưu trong bộ nhớ theo từng phiên và được dọn khi
phiên đóng.

\section{Cài đặt đối chiếu vị trí}

Khi sinh viên tự điểm danh, ứng dụng gửi kèm toạ độ, sai số định vị và cờ báo
vị trí giả lập do hệ điều hành cung cấp. Mã nguồn~\ref{lst:vitri} là phần
kiểm tra tương ứng.

\lstinputlisting[language=Python, caption={Đối chiếu vị trí khi sinh viên tự
điểm danh (đã lược chú thích).}, label={lst:vitri}]{code/doi-chieu-vi-tri.py}

Hàm không tự tìm địa điểm mà nhận vào vùng đối chiếu \code{fence} đã được
chọn sẵn theo thứ tự từ hẹp tới rộng: toạ độ riêng của phòng học, khu vực
chứa phòng, địa điểm gắn với lịch học, rồi địa điểm mặc định. Khi lấy toạ độ
riêng của phòng, bán kính được nâng lên tối thiểu \FbanKinhPhong\,m. Tách
riêng bước chọn vùng giúp hàm kiểm tra chỉ còn một việc là so khoảng cách,
và dễ kiểm thử với từng loại vùng.

Điểm cần lưu ý là cách xử lý sai số. Định vị trong nhà thường lệch hàng chục
mét \cite{Zandbergen2011Positional}, nên nếu so thẳng khoảng cách với bán kính thì rất nhiều sinh viên đang
ngồi trong lớp vẫn bị từ chối. Hệ thống trừ sai số khỏi khoảng cách trước khi
so sánh, đồng thời từ chối hẳn những báo cáo có sai số lớn hơn
\FsaiSoGps{}\,m vì khi đó toạ độ không còn giá trị đối chiếu.

\begin{luuy}
Toạ độ do thiết bị của sinh viên gửi lên nên về nguyên tắc vẫn có thể bị giả
mạo bằng ứng dụng chuyên dụng. Hệ thống chỉ coi đây là một lớp lọc bổ sung.
Quyền xác nhận cuối cùng vẫn thuộc về giảng viên đứng lớp.
\end{luuy}

\section{Cài đặt phân quyền theo phạm vi}

Phạm vi dữ liệu được cài đặt dưới dạng các hàm nhận vào một truy vấn và trả
về truy vấn đã bị thu hẹp, như Mã nguồn~\ref{lst:phamvi}.

\lstinputlisting[language=Python, caption={Thu hẹp truy vấn danh sách sinh
viên theo vai trò người đang đăng nhập.}, label={lst:phamvi}]{code/loc-pham-vi-sinh-vien.py}

Cách làm này có ba cái lợi. Thứ nhất, mọi endpoint liệt kê sinh viên đều đi
qua đúng một hàm, nên không thể quên áp phạm vi ở một chỗ nào đó. Thứ hai,
điều kiện được đẩy xuống câu lệnh SQL nên không có chuyện tải thừa dữ liệu
rồi lọc trên máy chủ ứng dụng. Thứ ba, khi một vai trò chưa được khai báo,
nhánh cuối cùng trả về truy vấn rỗng, tức mặc định là \term{từ chối} chứ
không phải \term{cho phép}, đúng nguyên tắc mặc định an toàn \en{(fail-safe defaults)} \cite{Saltzer1975Protection}, nhất là khi hệ thống mở rộng.

Hai vai trò cấp Khoa cùng đi qua một nhánh \code{FACULTY\_ROLES}: chúng thấy
cùng một phạm vi dữ liệu và chỉ khác nhau ở quyền chức năng. Văn phòng Khoa
được gọi các endpoint sửa lớp học phần, lịch học, ghi danh, nhưng mỗi lần sửa
đều kiểm tra lớp đó thuộc Khoa mình, còn lớp tạo mới thì tự gán vào Khoa của
người tạo. Lãnh đạo Khoa không có trong danh sách vai trò được sửa, nhưng là
vai trò được duyệt nghỉ có phép và xem nhật ký.

\section{Cài đặt đồng ý sinh trắc học và nhật ký hệ thống}

Nội dung đồng ý và số phiên bản \code{\FphienBanDongY} nằm trong một mô-đun
riêng ở phía máy chủ; ứng dụng chỉ tải về và hiển thị. Điều kiện đã đồng ý
được kiểm tra ở ba nơi độc lập: endpoint đăng ký khuôn mặt, câu truy vấn chọn
ứng viên khi quét camera (Mã nguồn~\ref{lst:nhandien}), và endpoint sinh viên
tự điểm danh. Kiểm tra ở cả ba nơi thay vì chỉ ở bước đăng ký để rút đồng ý
có hiệu lực tức thì, không phụ thuộc vào việc dữ liệu khuôn mặt đã được xoá
hết hay chưa.

Nhật ký hệ thống gồm hai phần. Hàm xác thực người dùng của mỗi yêu cầu gắn
thông tin người đang thao tác và địa chỉ IP vào phiên làm việc với cơ sở dữ
liệu. Sau đó một hàm lắng nghe sự kiện \code{after\_flush} của SQLAlchemy
duyệt các đối tượng vừa thêm, sửa, xoá, so lịch sử từng cột để lấy phần thay
đổi, rồi thêm dòng nhật ký vào cùng giao dịch. Những sự kiện không làm đổi dữ
liệu như đăng nhập hay xuất báo cáo thì được ghi tường minh tại endpoint
tương ứng.

\section{Tham số cấu hình}

Toàn bộ ngưỡng và quy tắc nghiệp vụ đều nằm ở tệp cấu hình, khai báo lại được
bằng biến môi trường, không có hằng số nào chôn trong mã nguồn.

{\scriptsize
\begin{xltabular}{\textwidth}{@{}>{\raggedright\arraybackslash\leavevmode}p{6.0cm} >{\centering\arraybackslash}p{1.6cm} X@{}}
\caption{Tham số cấu hình nghiệp vụ và giá trị mặc định.}
\label{tab:cauhinh}\\
\toprule
Tham số & Mặc định & Ý nghĩa \\
\midrule
\endfirsthead
\toprule
Tham số & Mặc định & Ý nghĩa \\
\midrule
\endhead
\bottomrule
\endfoot
\BangCauHinh
\end{xltabular}}

{\footnotesize
\begin{xltabular}{\textwidth}{@{}>{\bfseries\raggedright\arraybackslash\leavevmode}p{2.2cm} >{\raggedright\arraybackslash\leavevmode}p{5.9cm} >{\raggedright\arraybackslash\leavevmode}X@{}}
\caption{Giải thích chi tiết Bảng~\ref{tab:cauhinh}: các nhóm tham số và tác động khi đổi giá trị.}
\label{tab:cauhinh-gt}\\
\toprule
Nhóm & Tham số & Đổi giá trị thì \\
\midrule
\endfirsthead
\toprule
Nhóm & Tham số & Đổi giá trị thì \\
\midrule
\endhead
\bottomrule
\endfoot
Nhận diện & \code{RECOGNITION\_THRESHOLD}, \code{DUPLICATE\_THRESHOLD} & Tăng ngưỡng nhận diện thì ít nhận nhầm người hơn nhưng dễ không nhận ra chính chủ; giảm ngưỡng trùng thì chặn đăng ký hai danh tính chặt hơn nhưng dễ báo nhầm hai người giống nhau \\
Chống giả mạo & \code{ANTISPOOF\_ENABLED}, \code{ANTISPOOF\_THRESHOLD} & Tăng ngưỡng thì chặn ảnh giả chặt hơn nhưng dễ chặn nhầm người thật khi thiếu sáng \\
Xác nhận & \code{CONFIRM\_FRAMES} & Tăng thì chắc chắn hơn nhưng ghi nhận chậm hơn \\
Giờ giấc phiên & \code{EARLY\_OPEN\_MINUTES}, \code{LATE\_AFTER\_MINUTES}, \code{SESSION\_GRACE\_MINUTES} & Theo quy định của nhà trường về giờ mở điểm danh, đi muộn và hạn chót \\
Vị trí & \code{SELF\_CHECKIN\_REQUIRE\_LOCATION}, \code{SELF\_CHECKIN\_MAX\_ACCURACY\_M}, \code{ROOM\_MIN\_RADIUS\_M} & Nới ra thì ít từ chối oan sinh viên ngồi trong phòng, nhưng dễ điểm danh từ xa hơn \\
Mã QR & \code{CHECKIN\_QR\_ROTATE\_SECONDS}, \code{SELF\_CHECKIN\_REQUIRE\_QR\_DEFAULT} & Chu kỳ ngắn thì ảnh chụp mã hết hạn nhanh hơn, nhưng sinh viên quét chậm dễ bị trượt \\
Đăng nhập & \code{ACCESS\_TOKEN\_EXPIRE\_MINUTES}, \code{REFRESH\_TOKEN\_EXPIRE\_DAYS}, các tham số \code{LOGIN\_...}, \code{OTP\_EXPIRE\_MINUTES} & Thời hạn ngắn và giới hạn chặt thì an toàn hơn nhưng người dùng phải đăng nhập lại nhiều hơn \\
Dữ liệu, thông báo & \code{EVIDENCE\_RETENTION\_DAYS}, \code{MEDIA\_URL\_EXPIRE\_MINUTES}, \code{ABSENCE\_WARN\_RATIO}, \code{NOTIFY\_EMAIL\_...}, \code{APP\_PUBLIC\_URL} & Thời gian giữ ảnh, hạn của đường dẫn ảnh, ngưỡng cảnh báo cấm thi, cách gửi email và địa chỉ gắn vào link \\
Cách đặt & -- & Đặt biến môi trường cùng tên hoặc ghi vào tệp \code{backend/.env}, rồi khởi động lại máy chủ \\
\end{xltabular}}

Việc tách bạch như vậy cho phép điều chỉnh hệ thống theo thực tế từng đơn vị
mà không phải sửa mã: một khoa muốn siết ngưỡng nhận diện, một cơ sở có phòng
học trải rộng muốn nới bán kính điểm danh, đều chỉ là đổi giá trị cấu hình.

\section{Ứng dụng đa nền tảng}

Ứng dụng Flutter dùng chung một mã nguồn cho ba nền tảng, gồm
\FsoManHinh{} màn hình chia theo \FsoNhomManHinh{} nhóm vai trò. Điều hướng
do \code{go\_router} đảm nhiệm và chặn ngay ở tầng định tuyến: một tài khoản
sinh viên gõ thẳng đường dẫn của trang quản trị sẽ bị đưa về trang chủ, không
phụ thuộc vào việc menu có hiển thị mục đó hay không. Menu của mỗi vai trò
được xếp theo ma trận ở Bảng~\ref{tab:phanquyen}: việc chính của nhóm đứng
đầu, tiếp đến dữ liệu chỉ xem, nên văn phòng Khoa mở ứng dụng là thấy lớp và
lịch, còn lãnh đạo Khoa thấy báo cáo, duyệt phép và nhật ký.

Bố cục được thiết kế theo cổng thông tin sinh viên của trường để người dùng
không phải làm quen lại. Trên màn hình rộng, thanh trên cùng có logo, ô tìm
nhanh chức năng, chuông thông báo và tài khoản; menu chia nhóm được giấu sau
một nút dính ở mép trái và chỉ trượt ra khi cần, nhờ vậy nội dung dùng trọn
chiều ngang. Trang chủ của mọi vai trò xếp thành ba tầng: thẻ thông tin cá
nhân bên cạnh các ô số, hàng lối tắt tới chức năng của vai trò đó, rồi các
thẻ biểu đồ và bảng. Trên điện thoại, ứng dụng dùng thanh điều hướng dưới
hoặc ngăn kéo, và với sinh viên thì thẻ lịch hôm nay được đưa lên ngay sau
thẻ thông tin, vì tự điểm danh là việc chính khi mở ứng dụng trên điện thoại.

\begin{table}[H]
\centering
\caption{Các nhóm màn hình chính của ứng dụng.}
\label{tab:manhinh}
\small
\begin{tabularx}{\textwidth}{@{}>{\bfseries}p{3.0cm} X@{}}
\toprule
Nhóm & Màn hình tiêu biểu \\
\midrule
Xác thực & Đăng nhập, quên mật khẩu bằng OTP, đặt lại và đổi mật khẩu, tài
  khoản kèm danh sách chức năng của vai trò \\
Sinh viên & Lịch tuần, đồng ý sinh trắc học, đăng ký khuôn mặt, tự điểm danh
  kèm quét mã QR, chuyên cần của tôi \\
Giảng viên & Phiên điểm danh, quét camera cả lớp, chiếu mã QR, kết quả phiên,
  báo cáo \\
Phụ huynh & Danh sách con, chuyên cần của từng con, hộp thư thông báo \\
Phòng ban & Thống kê tổng hợp theo Khoa và theo lớp, nhật ký của Khoa \\
Quản trị & Sinh viên, giảng viên, lớp học phần, lịch, phòng học, địa điểm,
  duyệt khuôn mặt, phụ huynh, đơn vị, cán bộ, đặt lại mật khẩu, nhật ký hệ
  thống \\
\bottomrule
\end{tabularx}
\end{table}

{\small
\begin{xltabular}{\textwidth}{@{}>{\bfseries\raggedright\arraybackslash\leavevmode}p{2.6cm} >{\raggedright\arraybackslash\leavevmode}p{4.6cm} >{\raggedright\arraybackslash\leavevmode}X@{}}
\caption{Giải thích chi tiết Bảng~\ref{tab:manhinh}: ai dùng nhóm màn hình nào và đường dẫn tiêu biểu.}
\label{tab:manhinh-gt}\\
\toprule
Nhóm & Ai dùng & Đường dẫn tiêu biểu \\
\midrule
\endfirsthead
\toprule
Nhóm & Ai dùng & Đường dẫn tiêu biểu \\
\midrule
\endhead
\bottomrule
\endfoot
Xác thực & Mọi vai trò & \code{/login}, \code{/forgot-password}, \code{/account} \\
Sinh viên & Sinh viên & \code{/}, \code{/schedule}, \code{/my-attendance}, \code{/face}, \code{/self-checkin/:id} \\
Giảng viên & Giảng viên & \code{/sessions}, \code{/sessions/:id/scan}, \code{/sessions/:id/qr}, \code{/reports} \\
Phụ huynh & Phụ huynh & \code{/}, \code{/notifications}, \code{/children/:id} \\
Phòng ban & Phòng Đào tạo, Phòng Công tác sinh viên, văn phòng và lãnh đạo Khoa, ban ngành & \code{/}, \code{/reports}, \code{/audit-logs} cho lãnh đạo Khoa \\
Quản trị & Quản trị viên; nhiều trang dùng chung với phòng ban & \code{/students}, \code{/classes}, \code{/rooms}, \code{/staff}, \code{/units}, \code{/audit-logs} \\
Ghi chú & -- & Đường dẫn ngoài quyền của vai trò bị chặn ở tầng định tuyến và đưa về trang chủ \\
\end{xltabular}}

Ứng dụng tự làm mới access token khi hết hạn: lớp gọi API bắt mã lỗi 401, gọi
endpoint làm mới rồi phát lại yêu cầu ban đầu, nên người dùng không bị đăng
xuất giữa chừng. Token lưu bằng kho lưu trữ an toàn của hệ điều hành.

Hình~\ref{fig:tudiemdanh} là luồng tự điểm danh chạy thật trên điện thoại,
Hình~\ref{fig:sinhvien} là các màn hình chính của sinh viên và phụ huynh,
Hình~\ref{fig:tongquan} là màn hình tổng quan của phòng ban.

\begin{figure}[H]
\centering
\includegraphics[width=0.35\textwidth]{kq-ma-qr.png}\hfill
\includegraphics[width=0.205\textwidth]{kq-tu-dd-1.png}\hfill
\includegraphics[width=0.205\textwidth]{kq-tu-dd-2.png}\hfill
\includegraphics[width=0.205\textwidth]{kq-tu-dd-3.png}
\caption{Giảng viên chiếu mã QR, sinh viên đối chiếu vị trí, chụp ảnh và nhận
kết quả.}
\label{fig:tudiemdanh}
\end{figure}

\begin{figure}[H]
\centering
\includegraphics[height=7.6cm]{kq-sv-trang-chu.png}\hfill
\includegraphics[height=7.6cm]{kq-sv-chuyen-can.png}\hfill
\includegraphics[height=7.6cm]{kq-sv-khuon-mat.png}\hfill
\includegraphics[height=7.6cm]{kq-ph-thong-bao.png}
\caption{Trang chủ, chuyên cần, dữ liệu khuôn mặt của sinh viên và thông báo
của phụ huynh.}
\label{fig:sinhvien}
\end{figure}

\begin{figure}[H]
\centering
\includegraphics[width=0.85\textwidth]{kq-tong-quan.png}
\caption{Màn hình tổng quan chuyên cần của Phòng Đào tạo.}
\label{fig:tongquan}
\end{figure}

\section{Triển khai}

Toàn bộ hệ thống dựng lại được bằng một lệnh \code{docker compose up
-{}-build}, gồm ba dịch vụ: cơ sở dữ liệu PostgreSQL có sẵn pgvector, máy chủ
ứng dụng tự chạy di trú cơ sở dữ liệu trước khi phục vụ, và nginx phục vụ bản
Flutter Web đồng thời chuyển tiếp các yêu cầu API.

Hai lưu ý khi triển khai thật. Thứ nhất, cổng của cơ sở dữ liệu và máy chủ
chỉ mở trên máy cục bộ, người dùng đi vào qua nginx. Thứ hai, trình duyệt chỉ
cho phép truy cập camera trên HTTPS hoặc \code{localhost}, nên bản triển khai
thật bắt buộc phải có chứng chỉ HTTPS.

% ---------------------------------------------------------------------
%  CHƯƠNG 5 — KẾT QUẢ, ĐÁNH GIÁ VÀ KẾT LUẬN
% ---------------------------------------------------------------------
\chapter{Kết quả, đánh giá và kết luận}
\label{ch:ketqua}

\begin{muctieu}
Chương này đối chiếu sản phẩm thu được với từng mục tiêu đặt ra ở
Chương~\ref{ch:modau}, trình bày kết quả kiểm thử tự động, cách đo và kết
quả đo hiệu năng, bàn luận về những điểm hệ
thống làm được và chưa làm được, rồi tổng kết đồ án, nêu đóng góp chính và
các hướng phát triển tiếp theo.
\end{muctieu}

\section{Đối chiếu với mục tiêu ban đầu}

\begin{table}[H]
\centering
\caption{Mức độ hoàn thành từng mục tiêu nêu ở Mục~\ref{sec:muctieu}.}
\label{tab:doichieumuctieu}
\small
\begin{tabularx}{\textwidth}{@{}>{\bfseries}c >{\raggedright\arraybackslash}p{2.2cm} X@{}}
\toprule
Mục tiêu & Kết quả & Bằng chứng \\
\midrule
MT1 & Đạt & Quy trình đăng ký và nhận diện chạy trọn vẹn trên ONNX Runtime
  bản CPU, không cần GPU. Thời gian xử lý đo được nêu ở
  Bảng~\ref{tab:hieunang} \\
MT2 & Đạt & Sáu lớp phòng vệ ở Bảng~\ref{tab:phongve} đều đã cài đặt và có
  kiểm thử tự động tương ứng \\
MT3 & Đạt & \FsoEndpoint{} endpoint trên \FsoBang{} bảng dữ liệu phủ hết
  vòng đời từ hồ sơ tới báo cáo, có nhật ký sửa điểm danh, nhật ký hệ thống,
  thông báo cho phụ huynh và xuất Excel \\
MT4 & Đạt & Ma trận vai trò và chức năng ở Bảng~\ref{tab:phanquyen}, gồm cả
  việc tách văn phòng Khoa với lãnh đạo Khoa, được cài đặt tập trung và kiểm
  chứng bằng nhóm kiểm thử phân quyền \\
MT5 & Đạt & Một mã nguồn Flutter với \FsoManHinh{} màn hình chạy trên Web,
  Android, iOS; \FsoTest{} ca kiểm thử tự động chạy trên PostgreSQL thật \\
\bottomrule
\end{tabularx}
\end{table}

{\small
\begin{xltabular}{\textwidth}{@{}>{\bfseries\raggedright\arraybackslash\leavevmode}p{2.0cm} >{\raggedright\arraybackslash\leavevmode}p{5.6cm} >{\raggedright\arraybackslash\leavevmode}X@{}}
\caption{Giải thích chi tiết Bảng~\ref{tab:doichieumuctieu}: mỗi mục tiêu được kiểm chứng ở đâu.}
\label{tab:doichieumuctieu-gt}\\
\toprule
Mục tiêu & Yêu cầu tóm tắt & Kiểm chứng tại \\
\midrule
\endfirsthead
\toprule
Mục tiêu & Yêu cầu tóm tắt & Kiểm chứng tại \\
\midrule
\endhead
\bottomrule
\endfoot
MT1 & Nhận diện trọn vẹn trên máy chỉ có CPU, đủ nhanh để dùng trong lớp & Bảng~\ref{tab:hieunang} và Mục~\ref{sec:hieunang}; khung đông người còn chậm, ghi ở phần hạn chế \\
MT2 & Chống ảnh in hoặc màn hình, chống nhận nhầm do một khung xấu, chống điểm danh từ xa & Bảng~\ref{tab:phongve} và các tệp kiểm thử tương ứng ở Bảng~\ref{tab:kiemthu}; hiệu quả chặn ngoài thực tế chưa đo \\
MT3 & Số hoá trọn nghiệp vụ chuyên cần & Bảng~\ref{tab:api} và Bảng~\ref{tab:luocdo} \\
MT4 & Phân quyền đúng mô hình tổ chức & Bảng~\ref{tab:phanquyen} và các tệp kiểm thử phân quyền \\
MT5 & Một ứng dụng cho ba nền tảng, có kiểm thử tự động & Bảng~\ref{tab:manhinh} và Bảng~\ref{tab:kiemthu} \\
Cột Kết quả & Đạt nghĩa là có bằng chứng kiểm chứng được nêu ở cột cuối & Những điểm còn thiếu ghi ở mục hạn chế của chương này \\
\end{xltabular}}

\section{Kết quả kiểm thử tự động}

Bộ kiểm thử viết bằng pytest gồm \FsoTest{} ca trong \FsoFileTest{} tệp, khoảng
\FsoDongTest{} dòng mã kiểm thử. Bảng~\ref{tab:kiemthu} liệt kê từng nhóm.

{\scriptsize
\begin{xltabular}{\textwidth}{@{}>{\raggedright\arraybackslash\leavevmode}p{5.6cm} >{\centering\arraybackslash}p{1.3cm} X@{}}
\caption{Các nhóm kiểm thử tự động của phía máy chủ.}
\label{tab:kiemthu}\\
\toprule
Tệp kiểm thử & Số ca & Nội dung kiểm \\
\midrule
\endfirsthead
\toprule
Tệp kiểm thử & Số ca & Nội dung kiểm \\
\midrule
\endhead
\bottomrule
\endfoot
\BangTest
\end{xltabular}}

{\footnotesize
\begin{xltabular}{\textwidth}{@{}>{\bfseries\raggedright\arraybackslash\leavevmode}p{2.6cm} >{\raggedright\arraybackslash\leavevmode}p{6.4cm} >{\raggedright\arraybackslash\leavevmode}X@{}}
\caption{Giải thích chi tiết Bảng~\ref{tab:kiemthu}: các tệp kiểm thử theo nhóm rủi ro.}
\label{tab:kiemthu-gt}\\
\toprule
Nhóm & Tệp & Rủi ro được chặn \\
\midrule
\endfirsthead
\toprule
Nhóm & Tệp & Rủi ro được chặn \\
\midrule
\endhead
\bottomrule
\endfoot
Xác thực, phiên đăng nhập & \code{test\_auth}, \code{test\_tokens\_and\_media}, \code{test\_password\_reset}, \code{test\_change\_password}, \code{test\_manager\_reset\_password} & Lộ tài khoản, token bị đánh cắp được dùng lại \\
Phân quyền & \code{test\_authorization}, \code{test\_unit\_roles}, \code{test\_admin\_management}, \code{test\_admin\_crud} & Xem hoặc sửa dữ liệu ngoài phạm vi của vai trò \\
Nhận diện, chống gian lận & \code{test\_antispoof}, \code{test\_confirm\_frames}, \code{test\_location\_checkin}, \code{test\_checkin\_qr}, \code{test\_consent}, \code{test\_face\_approval} & Điểm danh hộ, xử lý dữ liệu sinh trắc khi chưa có đồng ý \\
Vòng đời điểm danh & \code{test\_open\_session\_rules}, \code{test\_session\_lifecycle}, \code{test\_session\_detail\_and\_records}, \code{test\_rooms}, \code{test\_week\_schedule} & Mở phiên sai giờ, xếp trùng phòng, thiếu bản ghi vắng \\
Báo cáo, thông báo, nhật ký & \code{test\_reports}, \code{test\_reports\_export\_and\_warnings}, \code{test\_my\_attendance\_detail}, \code{test\_parent\_notifications}, \code{test\_audit\_log}, \code{test\_evidence\_cleanup} & Số liệu sai, quên báo phụ huynh, thiếu dấu vết thao tác \\
Lỗi đã gặp & \code{test\_flow\_fixes} & Lỗi đã sửa quay trở lại \\
Cách đếm & -- & Số ca là số hàm kiểm thử trong tệp, đếm tự động; kiểm thử của ứng dụng Flutter chạy riêng bằng \code{flutter test}, không tính vào bảng \\
\end{xltabular}}

Ba quyết định về cách kiểm thử đáng được nêu rõ:

\begin{itemize}[leftmargin=1.6em, itemsep=3pt]
  \item \textbf{Chạy trên PostgreSQL thật, không dùng cơ sở dữ liệu nhúng.}
        Kiểu \code{Vector(\FchieuVector)} và toán tử khoảng cách của pgvector
        không có tương đương trong SQLite, nên nếu kiểm thử bằng cơ sở dữ
        liệu nhúng thì đúng phần quan trọng nhất lại không được kiểm.
  \item \textbf{Mỗi ca chạy trong một giao dịch rồi quay lui.} Nhờ vậy các ca
        không ảnh hưởng lẫn nhau mà không phải dựng lại lược đồ mỗi lần.
  \item \textbf{Giảm số vòng băm mật khẩu khi kiểm thử.} Bcrypt cố tình chậm \cite{Provos1999Bcrypt};
        hạ số vòng trong môi trường kiểm thử giúp bộ kiểm thử chạy nhanh mà
        không đổi logic nghiệp vụ.
\end{itemize}

Kết quả lần chạy gần nhất: toàn bộ \FsoTest{} ca đều đạt, thời gian chạy
39,56 giây (\code{278 passed in 39.56s}), thực hiện ngày 12/11/2026 trên
PostgreSQL 16 chạy trong Docker.

\section{Đánh giá hiệu năng}
\label{sec:hieunang}

\subsection{Cách đo}

Hiệu năng được đo ở ba điểm tách biệt: thời gian phát hiện và trích đặc trưng
một khung hình, thời gian truy vấn so khớp trong cơ sở dữ liệu, và thời gian
xử lý trọn vẹn một yêu cầu nhận diện tính từ lúc máy chủ nhận khung hình. Mỗi
phép đo lặp lại \HnSoLan{} lần, bỏ lần chạy đầu tiên vì lần đó còn phải nạp mô
hình, rồi lấy trung vị thay cho trung bình để giảm ảnh hưởng của các lần chạy
bất thường.

Phép đo do công cụ \code{backend/tools/do\_hieu\_nang.py} thực hiện ngày
\HnNgayDo{} trên máy thử nghiệm ở Bảng~\ref{tab:moitruong}. Vì chưa có ảnh
chụp bằng thiết bị sẽ dùng khi bảo vệ, phép đo dùng ảnh mẫu đi kèm thư viện
InsightFace: ảnh nhóm \code{t1.jpg} kích thước 1280$\times$886 có \HnSoMat{}
khuôn mặt, và một khung 640$\times$480 cắt từ chính ảnh đó, chứa đúng một khuôn
mặt. Phép truy vấn chạy trên cơ sở dữ liệu thử nghiệm có \HnSoVector{} vector
khuôn mặt. Công cụ ghi thẳng kết quả vào bảng dưới đây và vào các con số trong
phần nhận xét, nên số trong báo cáo là số máy in ra chứ không phải số chép tay.

\begin{table}[H]
\centering
\caption{Kết quả đo hiệu năng trên máy thử nghiệm.}
\label{tab:hieunang}
\small
\begin{tabularx}{\textwidth}{@{}X >{\centering\arraybackslash}p{3.0cm} >{\centering\arraybackslash}p{3.0cm}@{}}
\toprule
Phép đo & Trung vị & Lớn nhất \\
\midrule
\BangHieuNang
\bottomrule
\end{tabularx}
\end{table}

{\small
\begin{xltabular}{\textwidth}{@{}>{\bfseries\raggedright\arraybackslash\leavevmode}p{2.8cm} >{\raggedright\arraybackslash\leavevmode}X >{\raggedright\arraybackslash\leavevmode}p{4.6cm}@{}}
\caption{Giải thích chi tiết Bảng~\ref{tab:hieunang}: mỗi phép đo gồm những gì.}
\label{tab:hieunang-gt}\\
\toprule
Phép đo & Gồm những bước & Đầu vào \\
\midrule
\endfirsthead
\toprule
Phép đo & Gồm những bước & Đầu vào \\
\midrule
\endhead
\bottomrule
\endfoot
Một khuôn mặt & Chạy cả bộ mô hình đang nạp: phát hiện, ba mô hình phụ, trích đặc trưng & Khung 640$\times$480 cắt từ ảnh mẫu, đúng một khuôn mặt \\
Nhiều khuôn mặt & Như trên, trên cả ảnh nhóm & Ảnh mẫu \code{t1.jpg} 1280$\times$886, \HnSoMat{} khuôn mặt \\
Truy vấn so khớp & Một câu SQL tìm vector gần nhất bằng toán tử cosine của pgvector & Cơ sở dữ liệu thử nghiệm có \HnSoVector{} vector; vector truy vấn ngẫu nhiên \\
Chống giả mạo & Cắt vùng mặt mở rộng 2,7 lần, đưa về 80$\times$80, chạy MiniFASNetV2 & Khuôn mặt trong khung một khuôn mặt \\
Trọn vẹn một yêu cầu & Hàm \code{recognize\_frame} mà endpoint quét camera gọi: giải mã ảnh, chạy mô hình, chống giả mạo, so khớp trong phạm vi lớp; không gồm thời gian truyền mạng & Khung một khuôn mặt, lớp học phần đầu tiên trong cơ sở dữ liệu \\
Trung vị, lớn nhất & Trung vị và lần chậm nhất trong \HnSoLan{} lần đo, đã bỏ lần đầu vì lần đó còn nạp mô hình & -- \\
\end{xltabular}}


\subsection{Nhận xét về hiệu năng}

Gần như toàn bộ thời gian nằm ở phần chạy mô hình. Phát hiện và trích đặc
trưng một khuôn mặt mất \HnMotMat{}\,ms, trong khi kiểm tra chống giả mạo chỉ
mất \HnChongGiaMao{}\,ms và truy vấn so khớp \HnTruyVan{}\,ms. Con số truy vấn
nhỏ một phần vì cơ sở dữ liệu thử nghiệm ít dữ liệu; khi chạy thật, mỗi lần
so khớp chỉ xét vector của sinh viên trong lớp đang điểm danh, nên số vector
phải so vẫn chỉ cỡ sĩ số một lớp dù toàn trường có bao nhiêu sinh viên.

Mỗi khuôn mặt thêm vào khung hình tốn thêm khoảng \HnMoiMatThem{}\,ms: khung
\HnSoMat{} khuôn mặt mất \HnNhieuMat{}\,ms. Màn quét camera của giảng viên chụp
2 giây một lần và chỉ gửi khung mới khi khung trước đã xử lý xong, nên khung
đông người làm giãn nhịp quét chứ không làm dồn yêu cầu. Theo cách tính này,
một khung chụp trọn một lớp bốn mươi người sẽ mất nhiều giây; đây là ước
lượng, chưa đo trên lớp thật.

Phép đo cũng chỉ ra một chỗ tối ưu được ngay. Hệ thống đang nạp cả năm mô hình
của bộ \code{buffalo\_l}, trong đó ba mô hình phụ không được dùng tới. Đo lại
trên cùng ảnh với cấu hình chỉ nạp phát hiện và trích đặc trưng: khung một
khuôn mặt còn \HnGonMotMat{}\,ms, giảm \HnGiamMotMat\%; khung \HnSoMat{} khuôn
mặt còn \HnGonNhieuMat{}\,ms, giảm \HnGiamNhieuMat\%. Vector đặc trưng thu được
giống hệt cấu hình đầy đủ. Thay đổi này chưa được áp dụng vào hệ thống.

\section{Bàn luận và bài học rút ra}

\textbf{Phần mô hình không phải phần khó nhất.} Việc ghép SCRFD và ArcFace
vào hệ thống mất ít công hơn nhiều so với việc thiết kế đúng các quy tắc bao
quanh nó: ai được xem ảnh bằng chứng, phiên đóng lúc nào, sửa điểm danh thì
lưu vết ra sao, sinh viên rút lại dữ liệu sinh trắc thì xoá những gì. Đây là
kết luận đáng chú ý nhất rút ra từ đồ án.

\textbf{Thu hẹp không gian tìm kiếm là quyết định hiệu quả nhất.} Chỉ so khớp
trong phạm vi lớp học phần vừa giảm chi phí tính toán, vừa cắt giảm mạnh xác
suất nhận nhầm, vì hai người ngẫu nhiên giống nhau trong một lớp bốn mươi
người hiếm hơn nhiều so với trong cơ sở dữ liệu toàn trường.

\textbf{Sai số GPS trong nhà là rào cản thực tế.} Trong quá trình thử, việc
đối chiếu vị trí phải nới nhiều lần mới dùng được: trừ sai số thiết bị báo lên
trước khi so với bán kính, chỉ từ chối khi sai số vượt \FsaiSoGps{}\,m, và
kể cả khi phòng học có toạ độ riêng thì bán kính cũng không nhỏ hơn
\FbanKinhPhong\,m. Vì GPS không bao giờ đủ chính xác để phân biệt hai phòng
cạnh nhau, những buổi cần chắc chắn được giao cho mã QR xoay vòng. Đây là ví dụ cho thấy một tính
năng đúng về lý thuyết vẫn có thể không dùng được nếu bỏ qua đặc tính của
thiết bị thật.

\textbf{Phân quyền phải theo cách tổ chức làm việc thật.} Thiết kế ban đầu
gộp mọi cán bộ Khoa vào một vai trò vừa nhập liệu vừa duyệt nghỉ có phép. Khi
đối chiếu với cách Khoa vận hành, vai trò này được tách thành văn phòng Khoa
và lãnh đạo Khoa. Nhờ phạm vi dữ liệu đã được cài đặt tập trung, việc tách
chỉ cần thêm một vai trò vào nhóm vai trò cấp Khoa và chỉnh danh sách vai trò
ở các endpoint sửa dữ liệu, không phải viết lại bộ lọc nào.

\textbf{Kiểm thử tự động không thay được việc dùng thử thật.} Phiên bản đầu
của mã QR vượt qua toàn bộ kiểm thử nhưng chỉ dùng được từ bên trong ứng dụng:
ảnh QR chứa một chuỗi mã, phải vào đúng mục quét mã mới đọc được, còn camera
thường của điện thoại không biết làm gì với chuỗi đó. Thiếu sót này không lộ
ra ở kiểm thử mà ở lúc chạy thử theo thao tác của người dùng, và là lý do mã
QR được đổi sang dạng đường link. Lớp chống giả mạo cũng gặp tình huống tương
tự: kiểm thử dùng mô hình giả nên đều đạt, trong khi mô hình thật lại chấm ảnh
người thật là giả mạo vì mã nguồn chia điểm ảnh cho 255, khác với cách mô hình
được huấn luyện \cite{Minivision2020Silent}. Lỗi chỉ lộ ra khi chạy với tệp mô
hình và ảnh thật.

\textbf{Ràng buộc nên đặt ở tầng thấp nhất có thể.} Những quy tắc như mỗi lịch
học chỉ có một phiên trong ngày được đặt thành ràng buộc của cơ sở dữ liệu, nên
hai yêu cầu gửi đồng thời cũng không phá được. Nếu chỉ kiểm tra bằng lệnh điều
kiện ở tầng ứng dụng thì lỗi kiểu này rất khó tái hiện và rất dễ lọt.

\textbf{Viết kiểm thử ngay khi cài đặt là lựa chọn tiết kiệm.} Nhóm kiểm thử về
các lỗi luồng nghiệp vụ đã phát hiện là bằng chứng rõ nhất: mỗi lỗi tìm ra đều
được cố định lại bằng một ca kiểm thử, nhờ đó những lần sửa sau không làm hỏng
lại thứ đã chạy đúng.

\section{Hạn chế của hệ thống}

\begin{enumerate}[leftmargin=1.8em, itemsep=3pt]
  \item \textbf{Chưa đo độ chính xác và khả năng chống giả mạo trên dữ liệu
        thật.} Hai phép đo này cần ảnh của nhiều người thật chụp ở nhiều điều
        kiện, có sự đồng ý của người được chụp, và cần thử trực tiếp bằng ảnh
        in, màn hình và video trước camera. Đồ án mới kiểm chứng phần này bằng
        kiểm thử tự động và chạy mô hình trên ảnh mẫu, nên chưa có con số về tỉ
        lệ nhận đúng hay tỉ lệ chặn giả mạo.
  \item \textbf{Khung hình đông người chậm trên CPU.} Mỗi khuôn mặt thêm vào
        tốn khoảng \HnMoiMatThem{}\,ms (Mục~\ref{sec:hieunang}), nên quét một
        khung có cả lớp sẽ mất nhiều giây.
  \item \textbf{Ngưỡng nhận diện cố định.} Hệ thống dùng một ngưỡng chung
        $\tau = \FnguongNhanDien$ cho mọi người, trong khi ngưỡng tối ưu có
        thể khác nhau tuỳ chất lượng ảnh đăng ký của từng sinh viên.
  \item \textbf{Chống giả mạo chỉ dựa trên ảnh tĩnh.} Mô hình bị động không
        chống được những kiểu tấn công tinh vi hơn như mặt nạ ba chiều \cite{Ramachandra2017Survey}.
  \item \textbf{Bộ đếm xác nhận đa khung hình không bền vững.} Bộ đếm nằm
        trong bộ nhớ tiến trình, nên khởi động lại máy chủ giữa giờ sẽ khiến
        các phiên đang mở phải nhận diện lại từ đầu.
  \item \textbf{Chưa chạy thử ở quy mô thật.} Hệ thống mới chạy trên máy phát
        triển và dữ liệu mẫu, chưa qua thử nghiệm với nhiều lớp học đồng thời
        trên hạ tầng của nhà trường.
\end{enumerate}

\section{Kết luận}

Đồ án đã xây dựng hoàn chỉnh FAMS, một hệ thống quản lý sinh viên tích hợp
điểm danh bằng nhận diện khuôn mặt, chạy được từ đầu tới cuối trên máy chủ
không có GPU. Sản phẩm gồm một dịch vụ với \FsoEndpoint{} endpoint trên
\FsoBang{} bảng dữ liệu, một ứng dụng Flutter \FsoManHinh{} màn hình dùng
chung cho Web, Android và iOS, phân quyền cho \FsoVaiTro{} vai trò theo đúng
cơ cấu tổ chức của trường, cùng \FsoTest{} ca kiểm thử tự động chạy trên
PostgreSQL thật.

Cả năm mục tiêu đặt ra ở Mục~\ref{sec:muctieu} đều đã đạt, với bằng chứng
đối chiếu ở Bảng~\ref{tab:doichieumuctieu}. Phần chưa trọn vẹn nằm ở khâu
kiểm định: hệ thống chưa được đo độ chính xác nhận diện và khả năng chống giả
mạo trên dữ liệu thật, và chưa chạy thử ở quy mô thật của nhà trường.

\section{Đóng góp của đồ án}

\begin{itemize}[leftmargin=1.6em, itemsep=3pt]
  \item \textbf{Một quy trình điểm danh nhiều lớp phòng vệ.} Thay vì coi
        nhận diện khuôn mặt là lời giải duy nhất, hệ thống xếp chồng sáu lớp
        kiểm tra: chất lượng ảnh, chống giả mạo, xác nhận đa khung hình, thu
        hẹp không gian so khớp, đối chiếu vị trí theo phòng học và mã QR xoay
        vòng.
  \item \textbf{Mô hình phân quyền hai chiều.} Quyền theo chức năng kết hợp
        với phạm vi dữ liệu theo Khoa, cài đặt tập trung dưới dạng bộ lọc
        truy vấn, cho phép thêm vai trò mới mà không phải sửa từng endpoint;
        mọi thao tác quản lý đều được ghi vào nhật ký hệ thống.
  \item \textbf{Thiết kế tôn trọng dữ liệu sinh trắc.} Đồng ý sinh trắc học
        là bắt buộc, có phiên bản nội dung và được kiểm tra ở mọi điểm xử lý;
        chỉ lưu vector đặc trưng chứ không lưu ảnh lâu dài, ảnh bằng chứng
        phát qua đường dẫn được ký điện tử và tự hết hạn, người dùng rút đồng
        ý là dữ liệu bị xoá ngay.
  \item \textbf{Bộ tài liệu tự đồng bộ với mã nguồn.} Mọi con số về quy mô
        trong báo cáo này được sinh tự động từ kho mã nguồn bằng công cụ mô
        tả ở mục C của phần Phụ lục, nên báo cáo không thể mô tả sai hệ
        thống mà không bị phát hiện.
\end{itemize}

\section{Hướng phát triển}

\begin{enumerate}[leftmargin=1.8em, itemsep=3pt]
  \item \textbf{Đo độ chính xác và khả năng chống giả mạo.} Thu thập ảnh có
        sự đồng ý của nhiều sinh viên ở nhiều điều kiện ánh sáng và góc chụp,
        tính tỉ lệ chấp nhận nhầm và từ chối nhầm \cite{Jain2004Biometric,
        ISO19795}, chọn ngưỡng nhận diện theo đường cong ROC
        \cite{Fawcett2006ROC} thay vì đặt theo kinh nghiệm; thử chống giả mạo
        bằng ảnh in, màn hình và video, ghi cả số lần chặn nhầm người thật.
  \item \textbf{Chỉ nạp mô hình cần dùng.} Giới hạn bộ \code{buffalo\_l} ở
        phát hiện và trích đặc trưng; phép đo ở Mục~\ref{sec:hieunang} cho
        thấy thời gian giảm \HnGiamMotMat\% với khung một khuôn mặt mà vector
        không đổi.
  \item \textbf{Ngưỡng thích nghi theo từng người.} Tính ngưỡng riêng dựa
        trên độ phân tán của các vector đã đăng ký của mỗi sinh viên, thay
        cho một ngưỡng chung.
  \item \textbf{Chống giả mạo mạnh hơn.} Bổ sung tín hiệu chiều sâu nếu
        thiết bị có cảm biến phù hợp, hoặc kết hợp thêm kiểm tra chủ động
        trong luồng sinh viên tự điểm danh.
  \item \textbf{Đưa bộ đếm xác nhận ra kho lưu trữ ngoài tiến trình.} Dùng
        một kho khoá và giá trị như Redis để phiên đang mở không bị mất
        trạng thái khi máy chủ khởi động lại, đồng thời chạy được nhiều tiến
        trình máy chủ song song.
  \item \textbf{Buộc đổi mật khẩu sau khi được đặt lại.} Hiện cán bộ đặt lại
        xong vẫn biết mật khẩu của người đó. Nên yêu cầu người dùng đổi sang
        mật khẩu riêng ngay ở lần đăng nhập kế tiếp.
  \item \textbf{Kết nối với hệ thống quản lý đào tạo của trường.} Đồng bộ
        danh sách lớp học phần và ghi danh thay vì nhập lại, đây là điều
        kiện bắt buộc nếu muốn triển khai thật.
  \item \textbf{Thông báo đẩy trên điện thoại.} Phụ huynh hiện nhận thông
        báo trong hộp thư của ứng dụng và qua email. Bước tiếp theo là thông
        báo đẩy qua dịch vụ của hệ điều hành, và gửi cảnh báo cho cả sinh viên
        ngay khi tỉ lệ vắng chạm ngưỡng \FtiLeCanhBao\%.
\end{enumerate}


% ---------------------------------------------------------------------
%  TÀI LIỆU THAM KHẢO VÀ PHỤ LỤC
% ---------------------------------------------------------------------
\backmatter

\cleardoublepage
\phantomsection
\addcontentsline{toc}{chapter}{Tài liệu tham khảo}
\bibliographystyle{plain}
\bibliography{refs}

% ---------------------------------------------------------------------
%  PHỤ LỤC
% ---------------------------------------------------------------------
\chapter{Phụ lục}
\markboth{Phụ lục}{}
% Phu luc nam sau \backmatter nen chuong khong duoc danh so va bo dem bang, hinh, doan ma
% khong ve 0: de nguyen thi bang o day mang so tiep theo cua Chuong 5. Danh so rieng PL.1,
% PL.2... (ca ban \theH... de lien ket trong PDF khong trung voi bang cua Chuong 5).
\setcounter{table}{0}\renewcommand{\thetable}{PL.\arabic{table}}\renewcommand{\theHtable}{PL.\arabic{table}}
\setcounter{figure}{0}\renewcommand{\thefigure}{PL.\arabic{figure}}\renewcommand{\theHfigure}{PL.\arabic{figure}}
\setcounter{lstlisting}{0}\renewcommand{\thelstlisting}{PL.\arabic{lstlisting}}\renewcommand{\theHlstlisting}{PL.\arabic{lstlisting}}

\section*{A. Hướng dẫn cài đặt và chạy hệ thống}
\phantomsection
\addcontentsline{toc}{section}{A. Hướng dẫn cài đặt và chạy hệ thống}

Cách nhanh nhất là dùng Docker, vì hệ thống cần PostgreSQL có phần mở rộng
pgvector.

\begin{lstlisting}[language=bash, caption={Dựng toàn bộ hệ thống bằng Docker Compose.}]
# 1. Get the source code
git clone <repository-url> FAMS
cd FAMS

# 2. Create the environment file, then edit POSTGRES_PASSWORD and SECRET_KEY
copy .env.example .env

# 3. Build and start database, backend and web frontend
docker compose up --build

# 4. Load sample data (units, staff accounts, classes, schedules)
docker compose exec backend python seed.py

# 5. Open the web application
#    http://localhost:5173
\end{lstlisting}

Chạy từng phần khi phát triển:

\begin{lstlisting}[language=bash, caption={Chạy máy chủ và ứng dụng khi phát triển.}]
# Backend: needs Python 3.10 and a running PostgreSQL with pgvector
cd backend
python -m venv venv && venv\Scripts\activate
pip install -r requirements.txt
copy .env.example .env
alembic upgrade head
python seed.py
uvicorn app.main:app --reload --port 8001
# API documentation: http://localhost:8001/docs

# Flutter app on the web
cd mobile
flutter pub get
flutter run -d chrome --web-port 5174 --dart-define=API_BASE_URL=http://localhost:8001

# Automated tests
cd backend && pytest
cd mobile && flutter test
\end{lstlisting}

\begin{luuy}
Lần chạy đầu tiên, thư viện InsightFace tự tải bộ mô hình \code{buffalo\_l}
nặng khoảng 300\,MB, nên cần mạng và một khoảng chờ. Mô hình chống giả mạo
phải tải riêng và đặt vào \code{backend/models/}; nếu thiếu tệp này, hệ thống
vẫn chạy nhưng bỏ qua bước kiểm tra chống giả mạo và ghi một dòng cảnh báo.
Trình duyệt chỉ cho phép dùng camera trên HTTPS hoặc \code{localhost}.
\end{luuy}

\section*{B. Tài khoản mẫu để chấm và chạy thử}
\phantomsection
\addcontentsline{toc}{section}{B. Tài khoản mẫu để chạy thử}

Các tài khoản dưới đây do \code{seed.py} tạo ra, dùng để hội đồng chạy thử
từng vai trò. Mật khẩu chỉ dành cho môi trường thử nghiệm.

\begin{table}[H]
\centering
\caption{Tài khoản mẫu theo từng vai trò.}
\label{tab:taikhoanmau}
\small
\begin{tabularx}{\textwidth}{@{}>{\raggedright\arraybackslash}p{6.0cm} >{\raggedright\arraybackslash}p{3.0cm} X@{}}
\toprule
Tài khoản & Mật khẩu & Vai trò \\
\midrule
\code{admin@fams.edu.vn}      & \code{admin123}  & Quản trị viên \\
\code{pdt01@fams.edu.vn}      & \code{pdt12345}  & Phòng Đào tạo \\
\code{ctsv01@fams.edu.vn}     & \code{ctsv1234}  & Phòng Công tác sinh viên \\
\code{khoa.cntt@fams.edu.vn}  & \code{khoa1234}  & Văn phòng Khoa CNTT \\
\code{ldk.cntt@fams.edu.vn}   & \code{ldk12345}  & Lãnh đạo Khoa CNTT \\
\code{ban01@fams.edu.vn}      & \code{ban12345}  & Ban ngành \\
\code{gv01@fams.edu.vn}       & \code{gv123}     & Giảng viên \\
\code{sv01@fams.edu.vn}       & \code{sv123}     & Sinh viên \\
\code{ph01@fams.edu.vn}       & \code{ph123}     & Phụ huynh của sv01 \\
\bottomrule
\end{tabularx}
\end{table}

{\small
\begin{xltabular}{\textwidth}{@{}>{\raggedright\arraybackslash\leavevmode}p{2.6cm} >{\raggedright\arraybackslash\leavevmode}X@{}}
\caption{Giải thích chi tiết Bảng~\ref{tab:taikhoanmau}: nên thử gì với từng tài khoản.}
\label{tab:taikhoanmau-gt}\\
\toprule
Tài khoản & Nên thử \\
\midrule
\endfirsthead
\toprule
Tài khoản & Nên thử \\
\midrule
\endhead
\bottomrule
\endfoot
\code{admin} & Menu đầy đủ, nhật ký hệ thống toàn trường, đặt lại mật khẩu cho mọi tài khoản \\
\code{pdt01} & Tạo lớp học phần, xếp lịch vào phòng học và thử xếp trùng phòng để thấy bị chặn \\
\code{ctsv01} & Duyệt khuôn mặt, ghi nhận đồng ý theo bản giấy, duyệt nghỉ có phép, đặt lại mật khẩu sinh viên và phụ huynh \\
\code{khoa.cntt} & Tạo lớp và lịch trong Khoa CNTT; không có nút duyệt nghỉ có phép \\
\code{ldk.cntt} & Duyệt nghỉ có phép, xem nhật ký của Khoa; không có nút sửa lớp \\
\code{ban01} & Chỉ thấy thống kê tổng hợp, không thấy tên hay ảnh sinh viên \\
\code{gv01} & Mở phiên, quét camera cả lớp, chiếu mã QR, xác nhận tay và sửa điểm danh kèm lý do \\
\code{sv01} & Đồng ý sinh trắc học, đăng ký khuôn mặt, tự điểm danh bằng link trong mã QR \\
\code{ph01} & Xem chuyên cần của con, nhận thông báo mỗi lần con được điểm danh \\
Lưu ý & Các mật khẩu trên chỉ dành cho môi trường thử; bản chạy công khai đã đổi toàn bộ mật khẩu \\
\end{xltabular}}

\section*{C. Công cụ sinh số liệu cho báo cáo}
\phantomsection
\addcontentsline{toc}{section}{C. Công cụ sinh số liệu cho báo cáo}

Mọi con số về quy mô hệ thống trong báo cáo này không gõ tay mà được quét
trực tiếp từ kho mã nguồn bằng một công cụ nhỏ viết bằng Python.

\begin{lstlisting}[language=bash, caption={Sinh lại số liệu của báo cáo.}]
# Regenerate every number, table and code excerpt used by the report
python baocao/tools/quet-du-an.py

# Check only: exits with code 1 when the report data is out of date
python baocao/tools/quet-du-an.py --kiem-tra

# Static check of the LaTeX sources before compiling
python baocao/tools/kiem-tra-tex.py

# Measure recognition timings and write them into the report table
cd backend && python tools/do_hieu_nang.py --anh-mot-mat <image> --ghi
\end{lstlisting}

Công cụ đọc mã nguồn và sinh ra bốn tệp bảng cùng một tệp định nghĩa các lệnh
số liệu, đặt trong thư mục \code{baocao/auto/}. Từ đó, số bảng dữ liệu, số
endpoint, số ca kiểm thử, các ngưỡng cấu hình, các đoạn mã in trong
Chương~\ref{ch:caidat} và mã commit của bản được mô tả đều lấy thẳng từ hệ
thống thật. Công cụ cũng so sánh với lần quét trước và ghi ra tệp
\code{auto/thay-doi.md} liệt kê những gì đã đổi, kèm gợi ý mục nào trong báo
cáo cần soát lại.

\begin{ghinho}
Báo cáo này mô tả bản mã nguồn ở nhánh \code{\FnhanhGit}, commit
\code{\FcommitCuoi}, ngày \FngayCommit. Số liệu được quét lần gần nhất ngày
\FngayQuet.
\end{ghinho}

\section*{D. Bảng tự kiểm trước khi nộp}
\phantomsection
\addcontentsline{toc}{section}{D. Bảng tự kiểm trước khi nộp}

\begin{tukiem}
  \item Đã chạy \code{python baocao/tools/quet-du-an.py} lần cuối, số liệu
        trong báo cáo khớp với mã nguồn nộp kèm.
  \item Đã tìm chuỗi \emph{CẦN BỔ SUNG} trong tệp PDF và không còn chỗ nào
        sót lại.
  \item Đã điền họ tên, mã số sinh viên và lớp ở đầu tệp \code{do-an.tex}.
  \item Đã chèn ảnh chụp màn hình thật vào \code{figs/} và mở phần hình ở
        Chương~\ref{ch:caidat}.
  \item Bảng hiệu năng đã đo lại bằng \code{backend/tools/do\_hieu\_nang.py} trên
        đúng máy dùng khi bảo vệ; cấu hình máy và ngày đo khớp với báo cáo.
  \item Hệ thống chạy được trên máy khác, theo đúng các bước ở Mục A.
  \item Đã chạy thử demo trên đúng môi trường sẽ dùng khi bảo vệ.
  \item Mọi hình, bảng, giải thuật và mã nguồn đều có chú thích, có nhãn và
        được nhắc tới trong thân bài.
  \item Mọi mục tiêu ở Chương~\ref{ch:modau} đều được đối chiếu ở
        Chương~\ref{ch:ketqua}.
  \item Mọi tài liệu trong \code{refs.bib} đều được trích dẫn ít nhất một lần.
  \item Đã khai báo phần nào có dùng công cụ trí tuệ nhân tạo hỗ trợ, và giải
        thích được toàn bộ hệ thống khi bị hỏi.
\end{tukiem}

\section*{E. Nhận xét của giảng viên hướng dẫn}
\phantomsection
\addcontentsline{toc}{section}{E. Nhận xét của giảng viên hướng dẫn}

\khungnhanxet[7.5cm]

\vspace{0.8em}
\noindent\begin{minipage}[t]{0.5\textwidth}
\color{daunavy}
{\color{daumaroon}\bfseries Điểm số}

\vspace{0.7em}
\dongchamnhan{Bằng số:}

\vspace{0.9em}
\dongchamnhan{Bằng chữ:}
\end{minipage}\hfill
\okyten{\dauGVHD}{Giảng viên hướng dẫn}


\end{document}
